\section{Baader--Snyder: Equational Unification and Commutative Constructions}
\label{sec:p12-n-c-035-05}

This source-specific presentation develops the equational objects,
transformation systems, commutative constructions and disjoint-theory
combination in \cite[\S\S3--6, retained chapter
pp.~469--519]{BaaderSnyder2001}. Equality modulo a theory
changes the solution question, not the finite-tree term model.
The syntactic interface in \cref{sec:p12-unification-extraction-adequacy}
is reused only where the theory is empty or an explicit syntactic subproblem
is constructed. Its single-MGU conclusion does not apply to all theories.

\subsection{Problems, comparison, and complete solution sets}

\begin{hbdefinition}{def:p12-eu-equational-problem}{The theory and its finite problem}
Let \(\mathcal F\) be a ranked signature and \(\mathcal V\) a countably
infinite variable alphabet. Terms in \(T(\mathcal F,\mathcal V)\) are
finite trees. A set of identities \(s\approx t\), denoted by \(E\),
generates the least congruence \(=_{E}\) containing those identities
and closed under substitution. The source also calls the defining set
\(E\) an equational theory. Write \(\operatorname{Sig}(E)\) for its
occurring function symbols and require
\(\operatorname{Sig}(E)\subseteq\mathcal F\).
An \(E\)-unification problem is a finite set
\[
 \Gamma=\{s_1=^{?}_{E}t_1,\ldots,s_n=^{?}_{E}t_n\}.
\]
An \(E\)-unifier is a finite-support substitution \(\sigma\) over
\(\mathcal F\) with \(s_i\sigma=_E t_i\sigma\) for every \(i\).
The set of all such substitutions is \(\mathcal U_E(\Gamma)\). Fresh variables
may occur in images. Substitution is simultaneous, and source composition
acts first by \(\sigma\): \(x\sigma\lambda=(x\sigma)\lambda\).
It neither substitutes an inserted image again in the same application
nor admits rational trees. \cite[Definition 3.1, pp.~469--470]{BaaderSnyder2001}
\index{unification!equational}\index{equational theory}
\end{hbdefinition}

\begin{hbdefinition}{def:p12-eu-generality}{Comparison on the problem variables}
For a set \(\mathcal X\) of variables the source's dotted instantiation
comparison is \(\sigma\mathrel{\underline{\lessdot}}_E^{\mathcal X}\theta\) when some
substitution \(\lambda\) satisfies
\[
                       x\theta=_E x\sigma\lambda\quad(x\in\mathcal X).
\]
Write \(\sigma\doteq_E^{\mathcal X}\theta\) when both comparisons
hold. For a problem always use \(\mathcal X=\operatorname{Var}(\Gamma)\),
not the whole alphabet. A complete set \(\mathcal C\) satisfies
\[
 \mathcal C\subseteq\mathcal U_E(\Gamma),\qquad
 \forall\theta\in\mathcal U_E(\Gamma)\ \exists\sigma\in\mathcal C\quad
                  \sigma\mathrel{\underline{\lessdot}}_E^{\mathcal X}\theta.
\]
It is minimal complete if distinct members are incomparable: for
\(\sigma,\tau\in\mathcal C\), comparison
\(\sigma\mathrel{\underline{\lessdot}}_E^{\mathcal X}\tau\)
implies equality \(\sigma=\tau\) of the chosen representatives.
A most general \(E\)-unifier is a unifier whose singleton is complete.
The empty set is minimal complete for an unsolvable problem.
\cite[Definitions 3.2--3.3, pp.~470--471]{BaaderSnyder2001}
\end{hbdefinition}

\begin{hbtheorem}{thm:p12-eu-minimal-sets}{proved=annex:p12.n-c-035.5.minimal-sets}{Minimal solution classes characterize minimal complete sets}
The comparison is a preorder. Quotient \(\mathcal U_E(\Gamma)\) by its symmetric
part and use the induced partial order. Let \(M\) be its minimal elements.
A minimal complete set exists exactly when every solution class lies above
a member of \(M\). In that case it contains exactly one representative
of every member of \(M\), and its cardinality is independent of those
representatives. Uniqueness is up to instantiation equivalence, not in
general up to variable renaming.
\cite[Lemmas 3.10--3.11, Theorem 3.12, Corollary 3.13, pp.~475--477]{BaaderSnyder2001}
\end{hbtheorem}

\begin{hbdefinition}{def:p12-eu-types}{Four types and three signature regimes}
For a solvable problem, a singleton minimal complete set gives type
\(1\), unitary; a finite nonsingleton one gives \(\omega\), finitary;
an infinite one gives \(\infty\), infinitary. Absence of any minimal
complete set gives type \(0\), zero. The source orders them
\(1<\omega<\infty<0\) and calls a theory's type the maximal problem
type in its specified signature regime. ``At most finitary'' means
unitary or finitary. Type zero does not mean unsolvable: the latter
has the empty complete set and creates no higher-type obstruction.

An elementary problem has \(\mathcal F=\operatorname{Sig}(E)\).
A problem with constants adds only free nullary symbols. A general
problem may add free symbols of any arity. ``Free'' here means absent
from the theory signature, not eligible to be substituted as a variable.
Images use \(\mathcal F\) and fresh variables. The regimes may have
different types and decision problems for one and the same theory.
\cite[Definitions 3.4 and 3.9, pp.~471 and 474]{BaaderSnyder2001}
\index{unification!unitary}\index{unification!finitary}
\index{unification!infinitary}\index{unification!type zero}
\end{hbdefinition}

\begin{hbexample}{def:p12-eu-comparison-warning}{Instantiation equivalence need not be a renaming}
For \(I=\{f(x,x)\approx x\}\), put
\(\sigma=\{x\mapsto y\}\) and \(\theta=\{x\mapsto f(y,z)\}\).
On \(\{x\}\), \(\sigma\) covers \(\theta\) by assigning
\(y\mapsto f(y,z)\), and \(\theta\) covers \(\sigma\) by assigning
\(z\mapsto y\). A variable permutation cannot identify the distinct
\(y,z\), so it cannot make \(f(y,z)\) into an \(I\)-equal single
variable. Thus the syntactic renaming theorem cannot be substituted for
this equivalence. Elementary ACU is unitary for comparison on the problem
variables, but not necessarily for comparison on the entire alphabet.
\cite[\S3.2.1, p.~473]{BaaderSnyder2001}
\end{hbexample}

\begin{hbmethod}{def:p12-eu-output-contracts}{Algorithms, decisions, and enumeration}
An \(E\)-unification algorithm terminates with a finite complete set on
every input. A minimal algorithm also returns a minimal set. A decision
procedure only decides whether a solution exists. A unification procedure
enumerates a possibly infinite complete set; the qualifier ``minimal''
concerns that set, not termination. Such enumeration need not stop on an
unsolvable input, or even on an input that happens to have a finite minimal
complete set. If comparison is decidable, a finite complete output can be
made minimal by deleting redundant members, retaining one representative
of each equivalence class before pruning strictly covered classes.
\cite[pp.~472--473]{BaaderSnyder2001}
\end{hbmethod}

\begin{hbtheorem}{thm:p12-eu-c-and-a}{proved=annex:p12.n-c-035.5.c-and-a}{A finite C-construction and an infinite A-example}
Let \(C=\{f(x,y)\approx f(y,x)\}\), with \(f\) binary. Replace
each side of each input equation by every \(C\)-equal term, solve each
resulting syntactic problem, and collect its successful MGUs. This gives
a finite complete set of \(C\)-unifiers. A minimal complete set is
obtained by retaining minimal representative classes.
For distinct free constants \(a,b\),
\(f(x,y)=^{?}_C f(a,b)\) has exactly the two incomparable answers
\(\{x\mapsto a,y\mapsto b\}\) and
\(\{x\mapsto b,y\mapsto a\}\) on \(\{x,y\}\).

Let \(A=\{f(f(x,y),z)\approx f(x,f(y,z))\}\), with free constant
\(a\). The problem \(f(a,x)=^{?}_A f(x,a)\) has the infinite minimal
complete set \(\{\sigma_n:n\geq1\}\), where \(x\sigma_n\) is the
\(f\)-product of \(n\) copies of \(a\). Finiteness of every single
term's \(A\)-class therefore does not make the C-construction a
complete associative algorithm.
\cite[Examples 3.6--3.7, pp.~471--472]{BaaderSnyder2001}
\end{hbtheorem}

\begin{proofboundary}{def:p12-eu-zero-example}{The source's type-zero example}
The source reports that for \(AI=A\cup\{f(x,x)\approx x\}\),
\(f(x,f(y,x))=^{?}_{AI}f(x,f(z,x))\) has no minimal complete set.
It is solvable (identify \(y,z\)), not an empty-solution example.
The chapter attributes the nonexistence proof to Baader's 1986 work,
and gives Schmidt-Schau\ss\'s independent example
\(f(z,f(a,f(x,f(a,z))))=^{?}_{AI}f(z,f(a,z))\) with a free constant.
An abstract descending chain is not either theory-specific proof.
\cite[Example 3.8, p.~472]{BaaderSnyder2001}
\end{proofboundary}

\subsection{Free-algebra and logical interfaces}

\begin{hbdefinition}{def:p12-eu-free-algebra}{The free-algebra problem}
Nontrivial \(E\) means that some model has more than one element,
equivalently distinct variables are not \(E\)-equal.
The quotient \(T(\mathcal F,\mathcal X)/{=_{E}}\) is the free
\(E\)-algebra on \(\mathcal X\). The category \(C_{\mathcal F}(E)\)
has finitely generated free algebras and their homomorphisms.
A problem with \(n\) equations determines parallel arrows
\[
 \sigma,\tau:T(\mathcal F,\{u_1,\ldots,u_n\})/{=_{E}}
       \longrightarrow T(\mathcal F,\mathcal X)/{=_{E}},
 \quad u_i\sigma=[s_i],\quad u_i\tau=[t_i].
\]
A unifier is an arrow \(\delta\) with \(\sigma\delta=\tau\delta\).
Most generality requires existence, not uniqueness, of a factor arrow
for every other unifier. \cite[\S3.3.2--3.3.3, pp.~477--482]{BaaderSnyder2001}
\end{hbdefinition}

\begin{hbtheorem}{thm:p12-eu-categorical-bridge}{proved=annex:p12.n-c-035.5.categorical-bridge}{Substitution solutions are precisely free-algebra factorizations}
Restriction to input generators and passage to \(E\)-classes turn
substitutions into these arrows. The input equation conditions are
equivalent to equality of composite arrows, and problem-variable
comparison is equivalent to factorization. This does not make a most
general unifier a coequalizer. For the empty theory and
\(f(x,y)=^{?}f(y,x)\), the map \(x\mapsto x,y\mapsto x\) into
the free algebra on \(\{x,z\}\) is most general, but a factor's
value on \(z\) is unforced. Restricting that target to \(\{x\}\)
removes this nonuniqueness in this example, not in every nonempty theory.
\end{hbtheorem}

\begin{hbdefinition}{def:p12-eu-lcr}{Linear constant restrictions}
A problem with linear constant restrictions consists of a problem with
constants and a strict linear order \(<\) on its variables and free
constants. A permitted unifier must satisfy
\[
          x<c\quad\Longrightarrow\quad c\text{ does not occur in }x\sigma.
\]
For \(x=^{?}f(c)\), the order \(x<c\) excludes the syntactic solution,
whereas \(c<x\) permits it. The corresponding positive sentences are
\(\exists x\forall y\,x=f(y)\) and
\(\forall y\exists x\,x=f(y)\), respectively. Validity-oriented
Skolemization gives \(x=^{?}f(h(x))\) and \(x=^{?}f(d)\),
respectively: these are not satisfiability-oriented Skolem replacements
of existential quantifiers. \cite[Definition 3.15 and example, pp.~479--480]{BaaderSnyder2001}
\end{hbdefinition}

\begin{proofboundary}{def:p12-eu-logical-boundary}{Logical reformulations are separate proof dependencies}
The source's Theorem 3.14 relates elementary unification to the positive
existential fragment, and unification with constants to the positive
universal--existential fragment, both in \(E\) and in its countably
generated free algebra. Theorem 3.16 relates general unification, linear
constant restrictions, and the full positive fragment; its proof is
expressly a sketch using combination. Their reductions retain conjunctions
of all input equations; disjunctive-normal-form conversion can be
exponential. These logical proof dependencies are not supplied by the
elementary semiring argument below.
\cite[Theorems 3.14 and 3.16, pp.~478--480]{BaaderSnyder2001}
\end{proofboundary}

\subsection{The source transformation systems G, B, and N}

\begin{hbdefinition}{def:p12-eu-gbn-states}{Solved bindings, equation contexts, and two different state types}
The source's syntactic system \(\mathcal U\) uses \(P;S\), where \(P\) is
a finite equation multiset and \(S\) is solved: each left variable
occurs exactly once in \(S\) and nowhere on any right side. Its
substitution \(\sigma_S\) is the simultaneous map given by those
bindings; it is idempotent. Write \([\sigma]\) for the solved bindings
of an idempotent substitution. The source symbol \(e[u]\) denotes
an equation with one distinguished subterm occurrence; it is not the
unit constant in ACU. All identity/rule variants in an inference use
variables fresh for the previous computation.

In \(\mathcal G\), a solution of \(P;S\) unifies \(P\) modulo \(E\)
and unifies \(S\) literally. \(\mathcal G\) has U's four nonfailure
transformations and Lazy Paramodulation, but not the syntactic failure
rules on \(P\). In \(\mathcal B\) and \(\mathcal N\), a state is \(P;C;S\):
\(P\) is the equation skeleton multiset; \(C\) is a set of
syntactic constraints; \(S\) is solved. A solution unifies \(P\)
modulo the represented theory and \(C,S\) literally. Constrain moves
\(e\sigma_S\), not bare \(e\), from \(P\) to \(C\).
Basic Elimination does not instantiate \(P\); standard Elimination
does. Hence the basic nonvariable-position test is on the uninstantiated
skeleton, excluding positions created inside previous substitution images.
\cite[\S2.2.3, Definition 4.1, \S4.3, pp.~454--455,488,495--499]{BaaderSnyder2001}
\end{hbdefinition}

The exact transformations are registered in
\hyperref[sheet:p12-equational-unification-narrowing]{the source-specific
G/B/N sheet}. The matching equation introduced by LP has a special
first-step restriction: when its left side is nonvariable, matching
roots are required and that equation must first undergo Decomposition.
It is not permission to run ordinary syntactic Clash or Occurs on
equational problem terms. For \(E=\{f(x)\approx x\}\), the
equational problem \(x=^{?}_E f(x)\) has identity as a unifier,
although the syntactic problem has an occurs failure.

\calculussheet{sheet:p12-equational-unification-narrowing}{Baader--Snyder Equational Transformation and Narrowing Systems}{Finite first-order terms; source systems G: $P;S$, B and N: $P;C;S$; $P$ is a multiset and $S$ is syntactically solved; $\sigma_S$ acts first in composition. G uses fresh identities from $E\cup E^{-1}$; directed narrowing uses fresh $l\to r$ from $R$. Equation contexts $e[u]$ distinguish an occurrence, not a unit constant. See the standing, mandatory matching decomposition and proof-domain qualifications in \cref{def:p12-eu-gbn-states,def:p12-eu-narrowing-standing}. Sources: \cite[pp.~455,488--503]{BaaderSnyder2001}.}
\hbrule{rule:p12-eu-u-clash}{U helper Symbol Clash}{\{f(\bar s)=^{?}g(\bar t)\}\uplus P';S}{\bot}{different ranked symbols; syntactic U helper only, never a G failure rule}
\hbrule{rule:p12-eu-u-occurs}{U helper Occurs Check}{\{x=^{?}t\}\uplus P';S}{\bot}{$x\ne t$ and $x\in\operatorname{Var}(t)$; syntactic U helper only, never a G failure rule}
\hbrule{rule:p12-eu-g-trivial}{G Trivial}{\{s=^{?}s\}\uplus P';S}{P';S}{\none}
\hbrule{rule:p12-eu-g-decomposition}{G Decomposition}{\{f(s_1,\ldots,s_n)=^{?}f(t_1,\ldots,t_n)\}\uplus P';S}{\{s_1=^{?}t_1,\ldots,s_n=^{?}t_n\}\uplus P';S}{same ranked symbol; $n=0$ is allowed}
\hbrule{rule:p12-eu-g-orient}{G Orient}{\{t=^{?}x\}\uplus P';S}{\{x=^{?}t\}\uplus P';S}{$x$ variable; $t$ not a variable}
\hbrule{rule:p12-eu-g-eliminate}{G Variable Elimination}{\{x=^{?}t\}\uplus P';S}{P'\delta;S\delta\cup\{x\approx t\}}{$x\notin\operatorname{Var}(t)$; $\delta=\{x\mapsto t\}$; preserve the solved-form invariant}
\hbrule{rule:p12-eu-g-lp}{G Lazy Paramodulation}{\{e[u]\}\uplus P;S}{\{l=^{?}u,e[r]\}\uplus P;S}{$l\approx r$ fresh in $E\cup E^{-1}$; $u$ not a variable; if $l$ nonvariable then its root equals that of $u$, and the new matching equation must undergo Decomposition before any other rule is applied to it}
\hbrule{rule:p12-eu-b-trivial}{B/N Trivial}{P;\{s=^{?}s\}\cup C';S}{P;C';S}{\none}
\hbrule{rule:p12-eu-b-decomposition}{B/N Decomposition}{P;\{f(s_1,\ldots,s_n)=^{?}f(t_1,\ldots,t_n)\}\cup C';S}{P;\{s_1=^{?}t_1,\ldots,s_n=^{?}t_n\}\cup C';S}{same ranked symbol; zero arity allowed}
\hbrule{rule:p12-eu-b-orient}{B/N Orient}{P;\{t=^{?}x\}\cup C';S}{P;\{x=^{?}t\}\cup C';S}{$x$ variable; $t$ not a variable}
\hbrule{rule:p12-eu-b-eliminate}{B Basic Variable Elimination}{P;\{x=^{?}t\}\cup C';S}{P;C'\delta;S\delta\cup\{x\approx t\}}{$x\notin\operatorname{Var}(t)$; $\delta=\{x\mapsto t\}$; substitution is not applied to skeleton $P$}
\hbrule{rule:p12-eu-n-eliminate}{N Variable Elimination}{P;\{x=^{?}t\}\cup C';S}{P\delta;C'\delta;S\delta\cup\{x\approx t\}}{$x\notin\operatorname{Var}(t)$; $\delta=\{x\mapsto t\}$; this replaces B's elimination, not an additional B rule}
\hbrule{rule:p12-eu-b-constrain}{B/N Constrain}{\{e\}\uplus P';C;S}{P';\{e\sigma_S\}\cup C;S}{the instantiated equation moves into syntactic constraints}
\hbrule{rule:p12-eu-b-lp}{B/N directed Lazy Paramodulation}{\{e[u]\}\uplus P;C;S}{\{e[r]\}\uplus P;\{l\sigma_S=^{?}u\sigma_S\}\cup C;S}{$l\to r$ fresh in $R$; $u$ nonvariable in the current skeleton; if $l$ nonvariable its root agrees with $u$, and first decompose the new matching equation before another rule acts on it; source's inverse-identity reading is distinguished in \cref{def:p12-eu-narrowing-standing}}
\hbrule{rule:p12-eu-b-clash}{B/N constraint Symbol Clash}{P;\{f(\bar s)=^{?}g(\bar t)\}\cup C';S}{\bot}{different ranked symbols; applies to $C$, not equational $P$}
\hbrule{rule:p12-eu-b-occurs}{B/N constraint Occurs Check}{P;\{x=^{?}t\}\cup C';S}{\bot}{$x\ne t$ and $x\in\operatorname{Var}(t)$; applies to syntactic $C$, not equational $P$}
\hbrule{rule:p12-eu-b-solve}{B Composite Solve}{\{e\}\uplus P';C;S}{P';C\nu;S\nu\cup[\nu]}{$\nu$ an idempotent syntactic MGU of $e\sigma_S$; represents Constrain followed by syntactic constraint solving}
\hbrule{rule:p12-eu-b-narrow}{B Composite directed Narrow}{\{e[u]\}\uplus P;C;S}{\{e[r]\}\uplus P;C\nu;S\nu\cup[\nu]}{$l\to r$ fresh in $R$; $u$ nonvariable; $\nu$ a syntactic MGU of $l\sigma_S,u\sigma_S$ obtained with the LP matching restriction; skeleton not instantiated in B}
\hbrule{rule:p12-eu-problem-decomposition}{Problem Decomposition}{\{f(s_1,\ldots,s_n)=^{?}f(t_1,\ldots,t_n)\}\uplus P';C;S}{\{s_1=^{?}t_1,\ldots,s_n=^{?}t_n\}\uplus P';C;S}{$f$ is not the root of any left side in $R$}
\hbrule{rule:p12-eu-variable-abstraction}{Variable Abstraction}{\{e[s]\}\uplus P;C;S}{\{e[x]\}\uplus P;\{x=^{?}s\}\cup C;S}{$x$ fresh; use reduced-term side conditions when this is part of an ordered narrowing strategy}
\hbrule{rule:p12-eu-propagation}{Propagation}{\{e[u]\}\uplus P';C;\{x\approx t[s]\}\cup S}{\{e[y]\}\uplus P';C;\{x\approx t[s],y\approx s\}\cup S}{$u\sigma_S=s$ is nonvariable; $y$ fresh; standing and source's refined measure stated in \cref{def:p12-eu-narrowing-refinements}}
\hbrule{rule:p12-eu-blocking}{Blocking}{P;C;S}{\bot}{some term in $S$ is $R$-reducible, or some term in $C$ is reducible below its root; this rejects a reduced-witness branch, not all possible nonreduced E-solutions}
\hbrule{rule:p12-eu-ordered-critical-pair}{Source ordered critical pair}{s\approx t\qquad l\approx r}{(s[r]_p)\mu\approx t\mu}{fresh disjoint variants, either equation direction; $s|_p$ nonvariable; $\mu$ a syntactic U MGU of $l,s|_p$; some $\tau$ has $s\mu\tau\not\preceq t\mu\tau$ and $s\mu\tau\not\preceq(s[r]_p)\mu\tau$; order standing in \cref{def:p12-eu-ordered-cr}}
\hbrule{rule:p12-eu-reduce-source}{Reduce (literal source binding output)}{\{e[u]\}\uplus P;C;S}{\{e[r\gamma]\}\uplus P;C;S\cup[\gamma]}{$l\to r$ fresh; $u$ nonvariable; match $l\gamma=u\sigma_S$ using only fresh rule variables; the retained bindings are not automatically reduced certificates; source p501's abbreviated $l\gamma=u$ is qualified in \cref{def:p12-eu-eager-strategy}}
\hbrule{rule:p12-eu-reduce-cleaned}{Reduce (isolated match bindings discarded)}{\{e[u]\}\uplus P;C;S}{\{e[r\gamma]\}\uplus P;C;S}{same fresh-only match; after fully substituting $r\gamma$, discard only its match-domain bindings, absent from all old data and original variables; separately attributed editorial variant}
\hbrule{rule:p12-eu-propagation-nonvariable}{Propagation (nonvariable skeleton correction)}{\{e[u]\}\uplus P';C;\{x\approx t[s]\}\cup S}{\{e[y]\}\uplus P';C;\{x\approx t[s],y\approx s\}\cup S}{$u$ itself nonvariable in P; $u\sigma_S=s$ nonvariable; $y$ fresh; solved-image reduced-witness invariant; literal source variant remains separately registered}
\hbrule{rule:p12-eu-n-solve}{N Composite Solve}{\{e\}\uplus P';C;S}{P'\nu;C\nu;S\nu\cup[\nu]}{$\nu$ an actual syntactic MGU of $e\sigma_S$; includes the N updates of the whole skeleton}
\hbrule{rule:p12-eu-n-narrow}{N Composite directed Narrow}{\{e[u]\}\uplus P;C;S}{(\{e[r]\}\uplus P)\nu;C\nu;S\nu\cup[\nu]}{$l\to r$ fresh; $u$ nonvariable; $\nu$ an actual U MGU of $l\sigma_S,u\sigma_S$ with first matching Decomposition; N updates the whole skeleton}
\hbrule{rule:p12-eu-b-redex-ordered-narrow}{B Redex Ordered Narrow}{\{e[u]\}\uplus P;C;S}{\{e'[r]\}\uplus P;C\nu\cup C_0;S\nu\cup[\nu]}{B Narrow followed by the finite earlier-context abstraction pass of \cref{def:p12-eu-redex-ordered-strategy}; $C_0$ stores its fresh literal equations; original prescribed positions visited once, no repeated variable aliases}
\hbrule{rule:p12-eu-n-redex-ordered-narrow}{N Redex Ordered Narrow}{\{e[u]\}\uplus P;C;S}{\{e'[r\nu]\}\uplus P\nu;C\nu\cup C_0;S\nu\cup[\nu]}{N Narrow followed by the finite earlier-context abstraction pass in the instantiated context; $C_0$ stores its fresh literal equations; grounding and position domain declared in \cref{def:p12-eu-redex-ordered-strategy}}
\hbrule{rule:p12-eu-b-narrow-lower-index}{B Narrow with lower-index restriction}{\{e[u]\}\uplus P;C;S}{\{e[r]\}\uplus P;C\nu;S\nu\cup[\nu]}{B Narrow for rule index j; $l\sigma_S\nu$ is not a literal instance of any lower-index left side; not an E-unifiability test}
\hbrule{rule:p12-eu-n-narrow-lower-index}{N Narrow with lower-index restriction}{\{e[u]\}\uplus P;C;S}{(\{e[r]\}\uplus P)\nu;C\nu;S\nu\cup[\nu]}{N Narrow for rule index j; same literal lower-index matching test; may be combined with finite redex-ordered abstraction}


\begin{hblemma}{thm:p12-eu-u-bridge}{proved=annex:p12.n-c-035.5.u-bridge}{The source U helper has the actual removing-binding interface}
For a finite syntactic equation multiset \(P\), source U starts
\(P;\varnothing\), removes each selected binding from P, updates
the other equations and old solved bindings, and appends the new
binding to S. Every complete run terminates, fails exactly for
unsolvable input, or returns an idempotent syntactic MGU \(\sigma_S\).
Its source state is not the retained-equation sheet state.
The factorization is \(\rho=\sigma_S\rho\) globally for every
input solution \(\rho\). This gives the exact helper used in the
critical-pair construction below, while preserving the already proved
literal interfaces. \cite[\S2.2.3--2.2.4, pp.~454--457]{BaaderSnyder2001}
\end{hblemma}

\begin{hbtheorem}{thm:p12-eu-gbn-soundness}{proved=annex:p12.n-c-035.5.gbn-soundness}{Successful G/B/N traces are equationally sound}
Every terminal \(\varnothing;S\) reached by \(\mathcal G\), or
\(\varnothing;\varnothing;S\) reached by \(\mathcal B\) or \(\mathcal N\),
returns an \(E\)-unifier \(\sigma_S\) of the original problem.
This covers G's identities in both directions, and the directed
narrowing rules in the sheet. Constraint Clash and Occurs are sound
only for the literal constraint component. Blocking has a different
reduced-witness standing and is not included in an unsatisfiability
claim about all solutions of a state.
\end{hbtheorem}

\begin{hbdefinition}{def:p12-eu-abstract-completion}{An explicit critical-pair closure for the completeness proof}
Fix a reduction ordering \(\succ\) that is well founded, monotone,
and total on ground terms. Add finitely many fresh constants if needed
to freeze parameters; they are absent from \(E\). Let \(F_0=E\).
For \(F_j\), take both directions of fresh, variable-disjoint identities
\(l_1\approx r_1,l_2\approx r_2\) with nonvariable left sides.
At every nonvariable position \(p\) of \(l_1\), run the source U
helper on \(l_2=^{?}l_1|_p\), select one successful run's actual
idempotent MGU \(\mu\), when one exists, and
add the identity
\[
                   r_1\mu\approx (l_1[r_2]_p)\mu.
\]
Include all these critical identities in \(F_{j+1}\), retaining
\(F_j\), and put \(F_\omega=\bigcup_{j\geq0}F_j\).
Let \(\operatorname{Gr}^{\succ}(F_\omega)\) contain every ground
instance in its decreasing direction. This is an explicit unrestricted
critical-pair closure used as an editorial proof companion. The source's
\(E_\omega\) in Definition 4.10 instead denotes its ordered
\(\operatorname{Cr}(E_j)\) closure. The two are not declared identical;
the unrestricted closure is sufficient to derive the original G result,
and its simulation eliminates every extra critical identity before output.
\cite[Definitions 4.4 and 4.10, Lemma 4.11, pp.~490,492--493]{BaaderSnyder2001}
\end{hbdefinition}

\begin{hbtheorem}{thm:p12-eu-cp-ground-closure}{proved=annex:p12.n-c-035.5.cp-ground-closure}{The constructed closure is ground convergent and equationally conservative}
For nontrivial \(E\), the closure in
\cref{def:p12-eu-abstract-completion} has exactly the same equational
congruence as \(E\). Its decreasing ground instances terminate and
are confluent on ground terms, including any chosen fresh parameter
constants. Distinct ground terms are equal modulo \(E\) exactly
when they have the same normal form under those instances.
No algorithm computing this possibly infinite closure or its normal
forms is asserted.
\end{hbtheorem}

\begin{hbdefinition}{def:p12-eu-ordered-cr}{The actual ordered critical-pair interface}
The cross-referred rewriting chapter defines ordered critical pairs as
follows. Take fresh, variable-disjoint variants of \(s\approx t\)
and \(l\approx r\), with either direction of either equation allowed.
Unify \(l\) with a nonvariable subterm \(s|_p\), using an actual
syntactic U output \(\mu\). Keep the critical equation
\((s[r]_p)\mu\approx t\mu\) only if \emph{some} substitution
\(\tau\) satisfies both tests
\[
 s\mu\tau\not\preceq t\mu\tau,
 \qquad s\mu\tau\not\preceq(s[r]_p)\mu\tau,
 \quad\text{where }a\preceq b\text{ means }a=b\text{ or }b\succ a.
\]
The witness need not be ground; failure of \(\preceq\) also permits
incomparability on terms with variables. These are tests on the common
source and its two branches, not merely on the two resulting sides.
Use this source interface for \(\operatorname{Cr}^{\succ}(E)\),
and put \(E_0=E\), \(E_{j+1}=E_j\cup\operatorname{Cr}^{\succ}(E_j)\),
\(E_\omega=\bigcup_jE_j\). A reduction ordering is well founded,
monotone and substitution-stable, and is total on the chosen ground
signature, including fresh parameter constants.
The rewriting source writes \(\operatorname{cp}(E)\) and \(>\);
the unification source writes \(\operatorname{Cr}(E)\) and
\(\succ\). Their roles are connected explicitly here.
\cite[Definition 4.10, p.~492]{BaaderSnyder2001}
\cite[Definitions 4.5 and 6.2, pp.~547,569]{DershowitzPlaisted2001Rewriting}
\end{hbdefinition}

\begin{hbtheorem}{thm:p12-eu-ordered-closure}{proved=annex:p12.n-c-035.5.ordered-closure}{Ordered closure and its genuine comparison with the unrestricted companion}
The source \(E_\omega\) in \cref{def:p12-eu-ordered-cr} has the
same equational congruence as E. Its decreasing ground instances
terminate and are ground confluent. This includes the trivial-theory
case; nontriviality is still required by the G theorem's stated scope.
For nontrivial E, its ground normal form of any term is the same as
the normal form under the separately defined unrestricted \(F_\omega\):
both are the unique least ground member of that term's E-class.
The equation sets and rewrite relations need not be identical.
No ability to decide the existential overlap tests or compute an
infinite closure is inferred from the proof.
For the ordered G proof route choose an ordering with the proper-subterm
property on open terms, such as LPO. Then the filters exclude inner
variable-left overlaps in a nontrivial theory; with common U selectors
each source closure level is included in its unrestricted companion.
The existing two-direction critical simulation applies to precisely
those admitted pairs. Ground normalization itself is proved directly
for the ordered closure, not borrowed from its companion.
\cite[Lemma 4.11, pp.~492--493]{BaaderSnyder2001}
\cite[Definition 6.2 and Theorem 6.3, p.~569]{DershowitzPlaisted2001Rewriting}
\end{hbtheorem}

\begin{hbexample}{def:p12-eu-ordered-filter-example}{An unrestricted overlap the ordered filter rejects}
Take \(E=\{f(x)\approx g(x),g(a)\approx b\}\), with a
lexicographic path ordering extending \(f>g>b>a\). The reverse
first identity overlapped with \(g(a)\approx b\) gives the
unrestricted pair \(b\approx f(a)\). Its source is \(g(a)\),
which is smaller than the branch \(f(a)\), so the first ordered
test fails for every witness. The forward ground reductions already
give \(f(a)\to g(a)\to b\). Thus rejecting that particular
overlap does not omit this equality, and equality of normal forms
does not establish identity of the two closure sets.
\end{hbexample}

\begin{hbtheorem}{thm:p12-eu-cp-elimination}{proved=annex:p12.n-c-035.5.cp-elimination}{Both critical-pair directions can be eliminated from a G proof}
A finite successful G trace using identities from \(F_\omega\)
can be replaced by a finite trace using only \(E\cup E^{-1}\),
whose returned substitution agrees on the initial problem variables.
This includes the two directions of a critical identity and root
overlaps. The source's Case Two must respect LP's mandatory first
Decomposition: the repaired trace performs that step immediately and
then paramodulates at the corresponding position in a child equation.
\cite[Lemma 4.12, pp.~493--494]{BaaderSnyder2001}
\end{hbtheorem}

\begin{hbtheorem}{thm:p12-eu-g-completeness}{proved=annex:p12.n-c-035.5.g-completeness}{The original arbitrary-theory G calculus is complete}
For nontrivial \(E\), every \(E\)-unifier \(\theta\) of an
initial finite problem \(P;\varnothing\) is covered, on
\(\mathcal X=\operatorname{Var}(P)\), by the \(\sigma_S\) of
some finite G trace ending \(\varnothing;S\).
The proof constructs and lifts actual innermost ground rewrites,
with a reduced substitution, then eliminates the constructed critical
identities. Thus it proves the G result in the original theory, not
unification only in a completed surrogate. Completeness is existential
over branches, not termination or success of every G run.
\cite[Theorem 4.2, pp.~489--494]{BaaderSnyder2001}
\end{hbtheorem}

\begin{hbdefinition}{def:p12-eu-narrowing-standing}{Directed narrowing and the domain of ground convergence}
In directed narrowing a numbered finite ordinary rewrite system
\(R=\{l_1\to r_1,\ldots,l_n\to r_n\}\) supplies fresh
\(l\to r\) variants, with nonvariable \(l\) and
\(\operatorname{Var}(r)\subseteq\operatorname{Var}(l)\).
The source's composite Narrow in \S4.4.1 explicitly uses this direction;
the basic LP display in \S4.3.1 refers back to G's restrictions without
restating whether inverse identities are included. An inverse-identity
LP calculus must not be silently identified with the directed composite
Narrow calculus. The convergence theorem also remains valid if these
directed rules are embedded in a larger inverse-identity search tree:
the constructed directed branch is still present, although the trees
and the weak-hypothesis counterexample are then different.

The completeness theorem below uses a reduction ordering and convergence
on finite terms, equivalently convergence on ground terms after every
finite extension by fresh parameter constants. Merely converging on
the closed terms of the original signature is weaker. The source's
Theorem 4.13 says ``ground convergent'' and grounds its witness in
the proof. Initial open image normalization already follows from
the source's reduction ordering, as proved in
\cref{thm:p12-eb-initial-normalization}; it requires no new-constant
reserve. The additional grounding domain is relevant to common
normal forms and progress, not to existence of an initial reduced
witness. Its stronger confluence requirement is made explicit here,
rather than assumed by the word ``ground''; the original distinctions
are retained in \cref{def:p12-eb-source-scope,def:p12-eb-progress-boundary}.
\cite[\S4.3 and Theorems 4.13--4.14, pp.~496--499]{BaaderSnyder2001}
\end{hbdefinition}

\subsection{Reduced Ground Instances, Basic Proofs, and the Hereditary Restrictions}
\label{sec:p12-eu-hereditary-source}

The source's Definitions 4.3--4.6 and Lemmas 4.7--4.9 supply a
specific reduced-instance route, distinct from the later arbitrary-theory
completion construction. Its objects and proof-length measure are retained
here, with explicit corrections to the individual-step descent argument.
\cite[pp.~489--495]{BaaderSnyder2001}

\begin{hbdefinition}{def:p12-eu-h-ground-instances}{The source's reduced substitutions and ground-instance sets}
Fix nontrivial E and a reduction ordering \(\succ\), total on the
chosen ground signature. An R-reduced substitution has every image
\(x\theta\), for \(x\in\operatorname{Dom}(\theta)\), R-normal.
Put
\[
\begin{aligned}
 \operatorname{Gr}(E)&=\{l\rho\approx r\rho:
     l\approx r\in E\cup E^{-1},\ l\rho,r\rho\text{ ground}\},\\
 R&=\operatorname{Gr}^{\succ}(E)
   =\{l\rho\to r\rho:l\rho\approx r\rho\in\operatorname{Gr}(E),
                                               l\rho\succ r\rho\},\\
 R_E&=\{l\rho\to r\rho\in R:
          \rho\text{ is R-reduced on the identity variables}\}.
\end{aligned}
\]
Use the generating substitution only on
\(\operatorname{Var}(l,r)\); arbitrary unrelated assignments are
not part of the instance's provenance. The source's ``entire set''
for reducedness is R. Thus \(R_E\) is a defined subset of R,
not an assumption that every ground instance is reduced.
\cite[Definitions 4.3--4.4, pp.~489--490]{BaaderSnyder2001}
\end{hbdefinition}

\begin{hblemma}{thm:p12-eu-h-minimal-instance}{proved=annex:p12.n-c-035.5.h-minimal-instance}{Minimal applicable pairs really have reduced generating images}
In a finite term admitting an R-redex, choose an applicable pair
\((l\rho,r\rho)\) lexicographically minimal under \(\succ\).
Every generating image is R-normal, including variables occurring
only on the identity's right side. Consequently that redex is an
\(R_E\)-redex. R and \(R_E\) have precisely the same normal
terms. The choice is an existence construction, not an effective
enumeration or a claimed finite branching bound.
\end{hblemma}

\begin{hbdefinition}{def:p12-eu-h-basic-proof}{Based steps and the evolving source skeleton}
A step from \(q\theta\) is based on q when its redex is at a
nonvariable position in q. With a fresh rule presentation
\(l\rho\to r\rho\), replace the corresponding skeleton
occurrence by r, retaining its variables uninstantiated. The ground
successor is \(q[r]\theta\rho\), and the next step must be based
on this new skeleton. A finite based sequence is reflexive or repeats
these operations. A basic rewrite proof consists of based sequences
from \(s\theta\) and \(t\theta\) to their common normal form.
Normality of the final common term is part of ``rewrite proof'';
a zero-step equality at a reducible term is not such a proof.
\cite[Definition 4.6 and preceding convention, pp.~488,490]{BaaderSnyder2001}
\end{hbdefinition}

\begin{hbtheorem}{thm:p12-eu-h-basic-normalization}{proved=annex:p12.n-c-035.5.h-basic-normalization}{The original reduced-instance normalization and basic-proof equivalence}
Suppose R is ground convergent and \(s\theta\) is ground, with
\(\theta\) \(R_E\)-reduced. Every R-normalization of it has the
same endpoint as a finite based normalization using only \(R_E\).
For ground \(s\theta,t\theta\) with R-reduced \(\theta\),
\(s\theta=_E t\theta\) holds exactly when a basic rewrite proof
to a common normal form exists. It may be chosen using \(R_E\).
\cite[Lemma 4.7 and Corollary 4.8, p.~491]{BaaderSnyder2001}
\end{hbtheorem}

\begin{hblemma}{thm:p12-eu-h-basic-transport}{proved=annex:p12.n-c-035.5.h-basic-transport}{Literal solving does not destroy a reduced basic witness}
Let a finite \(R_E\)-basic proof be witnessed by reduced ground
\(\alpha\). If a finite syntactic substitution \(\mu\) satisfies
\(\mu\alpha=\alpha\) on its relevant variables, applying \(\mu\)
to the proof's initial skeletons preserves a basic proof with exactly
the same ground steps. Future identity variants are chosen fresh for
\(\mu\). In particular the minimum basic-proof length cannot
increase under a finite literal matching solve preserving that witness.
\end{hblemma}

\begin{hbdefinition}{def:p12-eu-h-source-measure}{The printed four-tuple and its explicit reduced-proof reading}
For \(P;S\) and a reduced witness \(\theta\), the source uses
\((m,n_1,n_2,n_3)\): m totals the lengths of minimal basic
normal-form rewrite proofs of the equations in \(P\theta\);
\(n_1\) counts distinct variables in equations whose instantiated
sides are literally the same normal term; \(n_2\) counts their
term-symbol occurrences; \(n_3\) counts their nonvariable-left,
variable-right orientations. Counts include equation multiplicity,
and may be zero; the order is lexicographic on natural numbers.

For the corrected reduced-instance proof, write \(m_{R_E}\) for
that minimum with all steps explicitly restricted to \(R_E\).
It exists by \cref{thm:p12-eu-h-basic-normalization}. The remaining
three coordinates are the source's same counts. The distinction
matters: a shortest unrestricted R-proof may erase a nonnormal image.
The induction step below is LP followed by its mandatory first
Decomposition and a full finite syntactic solve of the introduced
matching equation, not a raw LP step alone.
\cite[Lemma 4.9 proof, pp.~491--492]{BaaderSnyder2001}
\end{hbdefinition}

\begin{hbexample}{def:p12-eu-h-measure-failures}{Three different claims about shortest proofs}
For \(E=\{f(a)\approx a\}\), input \(f(a)=^{?}a\) has
minimum normal-form proof length one. Raw LP creates
\(a=^{?}a\) and \(f(a)=^{?}f(a)\); the latter needs two
steps to its common normal form, so m increases to two.

For \(E=\{f(x)\approx x,f(f(x))\approx x\}\), input
\(f(f(a))=^{?}a\) has a one-step reduced-instance proof. LP
with the second identity followed only by its first Decomposition
leaves \(f(x)=^{?}f(a)\), witnessed by \(x=a\); its two
normalization legs still total two. Full literal matching solving
is needed before asserting the source tuple's strict decrease.

Finally take \(E=\{f(x)\approx a,g(b)\approx c\}\) and
\(f(g(b))=^{?}a\). Its one-step unrestricted basic proof uses
\(x\mapsto g(b)\), which is not reduced. A reduced-instance
proof instead reduces \(g(b)\) to c and then \(f(c)\) to a.
This is why existence of a basic proof does not make its unrestricted
shortest version a reduced-instance proof.
\end{hbexample}

\begin{hbtheorem}{thm:p12-eu-h-minimum-blocks}{proved=annex:p12.n-c-035.5.h-minimum-blocks}{Lemma 4.9 with the source proof-length objects and a justified finite block}
Let nontrivial E have ground-convergent R in the declared parameter
signature, and let reduced \(\theta\) E-solve
\(P;\varnothing\) with \(P\theta\) ground. For open images,
first freeze their finite parameter set and require these hypotheses
in that fresh-constant extension, as in
\cref{def:p12-eu-narrowing-standing}. There is a finite
G trace to \(\varnothing;S\), with
\(x\theta=x\sigma_S\lambda\) on
\(\mathcal X=\operatorname{Var}(P)\).
Each selected normal-equation U step decreases
\((m_{R_E},n_1,n_2,n_3)\), and each selected reduced-instance
LP/full-matching block strictly decreases its first coordinate.
Fresh variables introduced by LP do not change old witness images.
This is a proof using the source's minimal basic-proof route with
the repairs stated above, not the earlier ground-term-multiset route.
\cite[Lemma 4.9, pp.~491--492]{BaaderSnyder2001}
\end{hbtheorem}

\begin{hbdefinition}{def:p12-eu-h-hereditary-tests}{The two source hereditary constraints}
Full nonvariable overlap compatibility requires that, at every position
where both l and u have a nonvariable symbol, their symbols and arities
agree. It does not require variables to have matching names, nor does
it replace the literal matching problem by E-equality.
The historical basic constraint records occurrence provenance:
constructor occurrences introduced by Variable Elimination inside
a substituted image are ineligible for subsequent LP, even though
they are physically present in the raw instantiated problem. Original
and freshly inserted identity-template constructor occurrences remain
eligible. Decomposition inherits the occurrence marks; syntactic
solving is not itself prohibited at marked occurrences.
These are the two restrictions described in the source's \S4.2.
\cite[p.~495]{BaaderSnyder2001}
\end{hbdefinition}

\begin{hbtheorem}{thm:p12-eu-h-convergent-hereditary}{proved=annex:p12.n-c-035.5.h-convergent-hereditary}{Both hereditary restrictions preserve the reduced-instance trace}
Under the exact ground-convergence/reduced-witness hypotheses of
\cref{thm:p12-eu-h-minimum-blocks}, its finite trace can satisfy
both hereditary constraints. Each LP match has a literal ground
witness, proving compatibility at every shared nonvariable position.
Every substituted copy retained in P has a normal ground value;
therefore the chosen next rewrite cannot be inside it. This establishes
actual occurrence eligibility, not a definition naming an unspecified
``basic lifting lemma.''
\end{hbtheorem}

\begin{proofboundary}{def:p12-eu-h-arbitrary-refinement-boundary}{The arbitrary-theory refinement is a separate remaining source argument}
The source's \S4.2 refers the arbitrary-theory refinements to
Dougherty--Johann 1992, Socher-Ambrosius 1994, and Moser 1993, and
explicitly does not develop their detailed proofs. The theorem just
proved is the source reduced-instance/ground-convergent regime; it
does not erase the arbitrary-theory refinement obligation.
The existing arbitrary-theory G theorem alone cannot discharge it:
critical-pair simulation may introduce function occurrences through
its unifier, whose historical eligibility differs from a primitive
critical-identity template. A proof must actually transport those
marks and the full-overlap tests through every nested simulation,
or reconstruct the identified refinement's own proof method.
\cite[\S4.2, p.~495]{BaaderSnyder2001}
\end{proofboundary}



\begin{hbtheorem}{thm:p12-eu-bn-completeness}{proved=annex:p12.n-c-035.5.bn-completeness}{Actual B and N lifting under the declared convergence hypothesis}
For \(R\) as in \cref{def:p12-eu-narrowing-standing}, every
\(R\)-equational solution \(\theta\) of \(P;\varnothing;\varnothing\)
is covered on \(\operatorname{Var}(P)\) by a finite directed B
trace, and by a finite directed N trace, ending
\(\varnothing;\varnothing;S\). In B the skeleton is left
uninstantiated at elimination; in N it is instantiated. The lifting
uses actual nonvariable skeleton positions because the witness's
images are reduced; this fact is proved in the annex.
The existential proof measure does not establish termination of the
whole nondeterministic search tree.
\end{hbtheorem}

\begin{hbexample}{def:p12-eu-ground-domain-failure}{Original-signature ground convergence is insufficient for directed narrowing}
Over the signature with constant \(a\), unary \(g\), and binary
\(f\), let
\(R=\{f(x,y)\to a,\ f(x,x)\to g(x),\ g(a)\to a\}\).
The positive integer interpretation
\(a=1,g(v)=v+1,f(v,w)=v+w+2\) strictly decreases for each rule,
so R terminates. Every closed term in this original signature reduces
to \(a\), by induction on its tree, so the system is ground
confluent there. Its equational theory is nontrivial: a two-element
algebra with both operations constantly equal to \(a\) is a model.
The critical peak at \(f(x,x)\) proves \(g(x)=_E a\), so identity
solves \(g(x)=^{?}_E a\).
Directed Narrow on that problem can only use \(g(a)\to a\),
forcing \(x=a\); no such answer covers identity, because
\(x=_E a\) fails in the two-element model. Adding a fresh constant
\(c\) exposes the unjoined peak \(a\leftarrow f(c,c)\to g(c)\).
This is a failing case for the weak domain reading in the directed
family. The concrete inverse-basic trace in
\cref{def:p12-en-inverse-branch} instead returns identity through
two permitted protected-root LP steps; its ordinary transport uses
that same inverse-identity family, not directed-forward Narrow.
\end{hbexample}

\begin{hbexample}{def:p12-eu-basic-skeleton-case}{A substituted position and a reduced witness are different things}
Let \(R=\{g(a)\to b\}\), and start with
\(P=\{h(x)=^{?}h(b),x=^{?}g(a)\}\), empty C and S.
Constrain on the second equation followed by B Elimination produces
\(P=\{h(x)=^{?}h(b)\}\), \(S=\{x\approx g(a)\}\).
B may not narrow inside the retained variable occurrence \(x\),
although N instantiates it to \(g(a)\) and permits that position.
The B branch has no reduced witness; Blocking detects its reducible
solved image. A complete B route instead first narrows \(g(a)\)
in the second skeleton equation to \(b\), then constrains and
eliminates \(x=b\). Constraining \(h(x)=h(b)\) now yields a
literal identity and returns \(x\mapsto b\), an equational unifier
of the original problem. Rejecting the first branch is not a proof
that the original problem lacks solutions.
\end{hbexample}

\subsection{Basic-to-Ordinary Narrowing: Exact Trace and MGU Transports}
\label{sec:p12-equational-narrowing-transport}

The source's Corollary~4.14 says that the proof of basic narrowing also
works for ordinary narrowing and that the same transformation sequence
can be used. Its composite Solve/Narrow presentation then suppresses
the intervening syntactic steps. The transports below make those two
dependencies explicit, without identifying the two state spaces or
borrowing a new convergence hypothesis for a stepwise equivalence.
\cite[\S\S4.3.1--4.4.1, pp.~496--499]{BaaderSnyder2001}

\begin{hbdefinition}{def:p12-en-variants}{The exact two state interfaces and LP direction}
Use the source states and six rules in
\cref{def:p12-eu-gbn-states,sheet:p12-equational-unification-narrowing}.
Write \(A=\sigma_S\) for the simultaneous substitution of solved S
and \(D_S=\operatorname{Dom}(A)\). In B, P is the retained skeleton;
in N it receives each elimination substitution. C contains literal
constraints, not equations modulo R. The comparison scope is always
the initial problem variables, retained externally even when an
equation deletion removes one from the current state.
Throughout the state and trace theorems, B and N mean these six
primitive transformations only. Reduce, Propagation, redex abstraction,
Blocking and other refinements are not implicitly included in their
reachable-state invariant.

Fix the same chosen family \(\mathcal L\) of fresh LP identities
in both systems. For the directed composite of source \S4.4.1 it is
the forward rules \(l\to r\in R\). The source's B display in
\S4.3.1 inherits G's restrictions, while G selects identities from
\(E\cup E^{-1}\); if that inverse-identity reading is used,
take that same family in both systems instead. The transport is
valid for either fixed family, not an identification of their trees.
In every LP use, u is nonvariable at the selected skeleton occurrence;
if l is nonvariable, roots match and its newly introduced matching
constraint first undergoes Decomposition. Variants are fresh for the
whole earlier computation, as in source footnote 11.
\cite[\S4.1.1, pp.~488--489; \S4.3.1, p.~496]{BaaderSnyder2001}
\end{hbdefinition}

\begin{hblemma}{thm:p12-en-state-invariant}{proved=annex:p12.en.state-invariant}{The actual reachable solved-state and global factor invariant}
In every reachable nonfailure B or N state, S is solved,
\(D_S\cap\operatorname{Var}(C)=\varnothing\), and A fixes its own
range. At eliminating \(x=^{?}t\) with \(\delta=\{x\mapsto t\}\),
x is outside \(D_S\), t contains neither x nor a member of \(D_S\),
and the new solved substitution is exactly \(A\delta\) on all
variables. In N, also \(D_S\cap\operatorname{Var}(P)=\varnothing\),
so \(P A=P\). Every literal S-solution \(\rho\) satisfies
\(A\rho=\rho\) globally. Thus at a terminal state \(\rho\)
itself is a factor for the returned answer, including initial
variables no longer occurring in the current P or C.
\end{hblemma}

\begin{hbdefinition}{def:p12-en-compile}{Compilation of a basic state and retained occurrences}
For a reachable basic state define its ordinary representation by
\[
         \operatorname{Inst}(P;C;S)=(P A;C;S).
\]
This is an equation-occurrence multiset, not a set obtained by
collapsing distinct equations. A nonvariable occurrence at address p
of a basic term remains at the same address of its A-instance, with
the same root and arity. Positions strictly inside a substituted
variable image may exist only in the latter; they receive no basic
address. Keep the first-Decomposition obligation on the same new
matching constraint in both representations.
\end{hbdefinition}

\begin{hbtheorem}{thm:p12-en-trace-transport}{proved=annex:p12.en.trace-transport}{Every basic trace has an actual ordinary trace with the identical answer}
Every finite B trace starting from \(P_0;\varnothing;\varnothing\)
has an N trace using the same chosen LP identities, such that at
each corresponding state \(P_N=P_B\sigma_S\) and C,S are literally
the same. Each of the six source rules and each matching-first
Decomposition is transported. If the basic trace terminates
\(\varnothing;\varnothing;S\), the ordinary trace terminates
with exactly that S and exactly that returned substitution.
No termination or confluence of R is required for this trace theorem.
The source's optional N Constrain variant that transfers bare e
instead of \(e A\) has the same effect on every reachable N state,
since \(P A=P\).
It is one-directional: an N step inside a substituted image need
not have a corresponding B step.
\end{hbtheorem}

\begin{hbmethod}{def:p12-en-source-composites}{Concrete U-MGU blocks and their accumulator interface}
For Solve take the literal equation \(Q=\{e A\}\). For directed
Narrow at a selected nonvariable u using fresh \(l\to r\), take
\(Q=\{l A=^{?}u A\}\). In the latter case first perform its
mandatory matching Decomposition, then run the actual removing-
binding U helper on the matching constraints. Let its idempotent
output be \(\eta\). It fixes variables outside \(\operatorname{Var}(Q)\)
and introduces no variable; in particular it fixes \(D_S\).
The new solved bindings are \(S\eta\cup[\eta]\), and their
substitution is \(A\eta\), not \(\eta A\).

In B the remaining/replaced skeleton is left uninstantiated; in N
it receives \(\eta\). The source's printed old-constraint context
\(C\eta\) may retain literal identities that a concrete set-based
helper run already deleted when a new constraint duplicated an old
one. Compare those presentations after Trivial deletion, never by
reclassifying C as equational. At the C-empty block boundaries
intended in \S4.4.1, this distinction is absent and the two composite
outputs are literal implementations of the registered rules.
\cite[\S2.2.3--2.2.4, pp.~454--457; \S4.4.1, p.~499]{BaaderSnyder2001}
\end{hbmethod}

\begin{hbtheorem}{thm:p12-en-mgu-blocks}{proved=annex:p12.en.mgu-blocks}{Finite composite blocks preserve the exact local factor condition}
These blocks terminate. A successful output has solved substitution
\(A\eta\). For every literal solution of old S and Q,
\(A\eta\rho=\rho\) globally. The output's literal S and old
constraint context describe exactly the old literal S,C,Q solutions,
with the stated harmless deletion of literal identities. A failing
match proves absence of that branch's literal witness, not absence
of all equational solutions of the original problem. Every successful
Solve or Narrow block is backwards equationally sound for its
selected problem equation.
\end{hbtheorem}

\begin{hblemma}{thm:p12-en-actual-lift}{proved=annex:p12.en.actual-lift}{A chosen innermost reduction lifts through an actual composite MGU}
Let \(\rho\) be a finite ground R-reduced witness for a reachable
B state, R-solving P and literally solving C,S. Suppose a selected term of P under
\(\rho\) has an innermost forward redex \(l\gamma\), where R
has ordinary rules (nonvariable left sides and no right-only variables).
Choose the least indexed rule matching that redex. Then its address
is nonvariable in the basic skeleton. The fresh matching substitution
\(\gamma\) extends \(\rho\) by normal proper-subterm images;
the literal matching constraint has that extension as a solution.
The actual U block therefore succeeds, with an explicit factor of
its \(\eta\), and preserves all old witness values. Under that
factor the new skeleton term is exactly the selected ground reduct.
The ordinary compiler transports the same block and factor. This
local lemma does not assume that every needed global witness is
groundable and normalizable in the original signature. Its groundness
and reduction tests are relative to the specified witness signature;
no ground confluence is needed for this one chosen-step construction.
\end{hblemma}

\begin{hbexample}{def:p12-en-transport-example}{Elimination changes the ordinary skeleton, not the basic occurrence map}
Let \(R=\{g(a)\to b\}\) and start with
\(\{h(x)=^{?}h(b),\ x=^{?}g(a)\}\). Constrain the second
equation and eliminate x. The B skeleton is
\(\{h(x)=^{?}h(b)\}\), the N skeleton is
\(\{h(g(a))=^{?}h(b)\}\), and both have S with
\(x\approx g(a)\). This is exactly \(P_N=P_B\sigma_S\).
N may narrow the g(a) child; B has a variable at its corresponding
address and may not use that occurrence. The trace theorem does not
transport this N-only step backwards or assert equality of search
trees. A different basic branch first rewrites the g(a) occurrence
while it is still in its original problem skeleton.
\end{hbexample}

\begin{hbexample}{def:p12-en-inverse-branch}{The weak-domain directed failure has a successful inverse-basic trace}
Use R and \(g(x)=^{?}a\) from
\cref{def:p12-eu-ground-domain-failure}, but now permit the inverse
identity of \(f(x,x)\to g(x)\). Choose fresh
\(g(v)\approx f(v,v)\) at the original g-root. The matching
constraint first decomposes to \(v=^{?}x\); Basic Elimination
sets v=x while the retained skeleton becomes \(f(v,v)=^{?}a\).
Its instantiated representation is \(f(x,x)=^{?}a\).
Next choose the forward identity \(f(p,q)\approx a\), fresh p,q.
Its matching constraint is \(f(p,q)=^{?}f(x,x)\); first
Decomposition and the literal solver give p=x,q=x. The remaining
skeleton is a=a, which Solve deletes. The terminal S is
\(\{v\approx x,p\approx x,q\approx x\}\), whose answer on
the initial scope \(\{x\}\) is identity. Both LP positions were
roots of actual basic skeletons. The trace compiler gives an N
trace in the same inverse-identity family with the identical answer.
It does not give a trace in the smaller directed-forward family.
Thus the directed weak-domain counterexample cannot be transferred
to the raw inverse-basic interpretation merely by calling both
procedures narrowing.
\end{hbexample}

\begin{hbtheorem}{thm:p12-en-terminal-normalization}{proved=annex:p12.en.terminal-normalization}{Terminal normalization preserves the answer class, not an active literal state}
For a terminal answer A and an ordinary R, if each image of A has
a finite reduction to a chosen normal form, normalize those images
to obtain \(\widehat A\), identity outside the same domain. Then
\(A=_R\widehat A\) pointwise, \(\widehat A\) is still idempotent,
and the two answers cover exactly the same targets modulo R on
the original scope. Soundness and completeness of any returned
family are preserved by this postprocessing. Termination of the
image reductions suffices for this assertion; confluence makes
the chosen forms unique. It does not license replacing a binding
in active S without transporting its literal constraints.
\end{hbtheorem}

\begin{hbexample}{def:p12-en-active-normalization-failure}{Equationally equal bindings can invalidate the literal active witness}
For \(R=\{f(a)\to a\}\), take the active state
\(P=\varnothing\), \(C=\{z=^{?}f(a)\}\),
\(S=\{x\approx f(a)\}\). Its literal witness gives
x=z=f(a). Normalizing only the x-binding to a removes that
witness, since the new S would require x=a literally. E-equality
between f(a) and a does not repair a literal S-equation. At a
terminal state with no C, normalized x=a is a sound representative
of the same equational answer class. These are different operations.
\end{hbexample}

\begin{hbtheorem}{thm:p12-en-infinite-search}{proved=annex:p12.en.infinite-search}{A finite convergent presentation can have an infinite basic and ordinary branch}
Let \(R=\{f(s(x))\to f(x),\ f(0)\to a\}\), with distinct
constants 0,a and unary f,s. R terminates and is confluent on all
finite terms. Starting from \(f(y)=^{?}a\), there are successful
source composite B and N traces returning \(y\mapsto s^k(0)\)
for every \(k\geq0\), and an infinite branch that repeatedly
chooses the first rule. Restricted to the initial scope \(\{y\}\),
these answers form an infinite minimal complete set for that problem,
not just infinitely many redundant traces. Thus neither finite convergent R nor finite
matching U blocks imply termination of the nondeterministic narrowing
search. The repeated B steps occur in fresh replacement skeletons,
not inside previously bound images.
\end{hbtheorem}

\begin{proofboundary}{def:p12-en-completeness-boundary}{The transported component and the remaining universal dependency}
The trace theorem supplies the actual source Corollary~4.14 transport:
every constructed successful basic witness trace is an ordinary trace
with the same answer and comparison factor. The MGU and lifting
lemmas supply the concrete \S4.4.1 composite blocks used on that
witness branch. They add no universal completeness premise and
do not replace the existing B lifting argument by a named result.

Initial open image normalization is supplied at the actual source
reduction-ordering scope by
\cref{thm:p12-eb-initial-normalization}; it needs neither grounding
nor an infinite/new-constant reserve and does not alter active S.
When one original ground term exists, the separate termination
argument in \cref{thm:p12-eb-ground-termination} also supplies open
termination under the unfiltered rewrite convention. Neither result
turns fixed-signature ground confluence into open confluence.

Example~\ref{def:p12-eu-ground-domain-failure} gives the directed
weak-domain failure; Theorem~\ref{thm:p12-eu-bn-completeness} keeps
its separately declared strengthened confluence domain. For raw
inverse-basic B, the still-required argument is global progress,
equational conversion to basic paths and protected occurrence history,
not existence of a normalized initial image witness. The actual
countercase in \cref{thm:p12-eb-progress-countercase} has normal
compatible witness images and a successful inverse branch, but its
first step increases M. The precise unclosed dependency is
\cref{def:p12-eb-progress-boundary}; the stepwise compiler alone
does not close it.

These trace and factor results do not identify Fay's historical
presentation with normalization-based algorithms or establish the
arbitrary-theory hereditary refinements. In particular a run that
rewrites active S modulo R is not the ordinary N state transport
proved here.
\cite[Theorem4.13, Corollary4.14 and \S4.4.1, pp.~497--499]{BaaderSnyder2001}
\end{proofboundary}


\subsection{The Original Term Domain and the Remaining Inverse-Basic Progress Step}
\label{sec:p12-equational-basic-completeness}

The earlier definitions matter to the scope of source Theorem~4.13.
They permit open terms and open target substitutions. Initial image
normalization can be justified without introducing new constants;
the separate obstacle is progress from R-equivalent open normal
endpoints. The results below isolate that obstacle and give an actual
inverse-basic trace that preserves normal witness images while
increasing the source's first complexity coordinate.
\cite[\S2.1, pp.~451--452; \S4.3, pp.~495--498]{BaaderSnyder2001}

\begin{hbdefinition}{def:p12-eb-source-scope}{The source's term, substitution and convergence domains}
In Baader--Snyder \S2.1, \(\mathcal F\) is finite or countably
infinite, \(\mathcal V\) is countably infinite, and terms belong to
\(T(\mathcal F,\mathcal V)\). Nullary symbols belong to
\(\mathcal F\); no infinite reserve of nullary symbols is required.
Substitutions map variables to those terms and have finite support.
Their composition acts left to right. Definition~3.1 permits arbitrary
such substitutions as equational unifiers, not only ground ones.
Definition~4.3 calls a substitution R-reduced when all its nonidentity
images are R-normal. A B-state solution R-unifies P but literally
unifies C,S. These are different equality requirements.
\cite[pp.~451--452,469--470,489,495]{BaaderSnyder2001}

The rewriting terminology explicitly imported by \S4.1 is defined
on all finite terms over function symbols, constants and variables.
A rewrite step includes every contextual substitution instance of
an R-rule; it has no distinct-endpoint proviso. A derivation may
therefore contain a step whose endpoints coincide. A reduction
ordering is well founded, monotone and fully substitution-invariant.
Ordinary convergence means termination and confluence on the entire
term domain; ground confluence restricts the latter condition to
variable-free terms of that vocabulary. The two definitions are
not identified. Baader--Snyder's \S4.3/Theorem4.13 expressly says
ground convergent and does not separately declare closure under
free-constant extensions.
\cite[\S2, pp.~541--543; Definitions4.3--4.5, pp.~546--547;
Definitions5.5--5.6, pp.~560--561]{DershowitzPlaisted2001Rewriting}
\cite[pp.~488,496--497]{BaaderSnyder2001}
\end{hbdefinition}

\begin{hbdefinition}{def:p12-eb-ground-map}{One-ground-term evaluation is a proof map, not an infinite-support source substitution}
When \(q\in T(\mathcal F,\varnothing)\), let \(H_q\) evaluate
every variable at q and preserve every function symbol. The map acts
on individual finite terms. It need not be a finite-support
substitution on the infinite alphabet. On each finite set of terms
it agrees with the finite substitution assigning their finitely
many occurring variables to q. On any one rewrite step, a finite
restriction to the chosen rule variables realizes its grounded
instance. This distinction also applies when an infinite sequence
uses infinitely many variable names overall.
\end{hbdefinition}

\begin{hblemma}{thm:p12-eb-ground-termination}{proved=annex:p12.eb.ground-termination}{Ground termination supplies open termination when one ground term exists}
Under the source's unfiltered rewrite relation, if the original
signature has at least one ground term and R terminates on all
its ground terms, then R terminates on all its open terms.
This requires neither distinct ground codes for parameters nor
new constants. It implies no open confluence. If equal-endpoint
rewrite instances were discarded, or if no ground term existed,
the stated implication would need additional hypotheses.
\end{hblemma}

\begin{hblemma}{thm:p12-eb-initial-normalization}{proved=annex:p12.eb.initial-normalization}{Initial open reduced witnesses exist at the actual ordering scope}
Suppose \(l\succ r\) for every R-rule under a reduction ordering
on the source term domain. Then R terminates on open terms;
each left side is nonvariable and
\(\operatorname{Var}(r)\subseteq\operatorname{Var}(l)\).
For every finite-support R-unifier \(\theta\) of an initial
\(P;\varnothing;\varnothing\), normalize its finitely many
nonidentity images by any finite R-reductions. The resulting
finite-support \(\theta_0\) is R-reduced, still R-unifies P,
and is pointwise R-equal to \(\theta\). Thus \(\theta_0\) and
\(\theta\) cover each other on the initial problem variables
with the identity factor. No idempotence of the target substitution,
grounding, fresh-constant reserve or uniqueness of image normal
forms is needed. This is normalization of an initial witness while
S is empty, not alteration of an active literal binding.
\end{hblemma}

\begin{hbtheorem}{thm:p12-eb-open-history-lift}{proved=annex:p12.eb.open-history-lift}{Open normal images, actual helper factors and protected basic history}
Consider the six primitive B transformations, with the source solved
state and actual U helper. Keep a finite-support witness \(\rho\)
R-reduced, R-solving P and literally solving C,S. Literal constraint
steps preserve that witness, including its globally retained initial
scope. If a term of \(P\rho\) has a chosen innermost forward R-redex,
its address is nonvariable in P. A fresh rule match gives normal
proper-subterm images, and the actual matching-first U block succeeds
with computed \(\eta\) and explicit factor
\(\eta\rho'=\rho'\), where \(\rho'\) preserves all old values.
The new solved substitution is \(\sigma_S\eta\), and
\(\sigma_S\eta\rho'=\rho'\) globally.

P's eligible function occurrences originate in the initial skeleton
or a fresh LP right-side template, never in Basic Elimination.
Each protected substitution image has a normal value under the
maintained witness; hence a selected forward redex cannot be inside
it. These assertions hold for open witnesses and need no common-
normal-form or ground-confluence assumption. They do not assert
that a forward redex exists whenever two open endpoints are R-equal.
Nor does forward normality rule out inverse expansions inside an
image. The latter are excluded by the basic addresses, not by a
false bidirectional-normality claim.
\end{hbtheorem}

\begin{hblemma}{thm:p12-eb-inverse-expansion}{proved=annex:p12.eb.inverse-expansion}{A normal eligible inverse match has a compatible normal-witness extension}
Under the same ordering and primitive-state standing, let u be a
nonvariable basic skeleton occurrence with \(u\rho\) R-normal.
Suppose a fresh reversed R-rule \(r\approx l\), from forward
\(l\to r\), has an actual match \(r\gamma=u\rho\).
Its matched variable images are normal; any variables in l absent
from r can be assigned fresh normal parameter variables. The
matching-first helper then supplies a finite raw inverse-basic
block preserving all old witness values and producing the actual
replacement \(l\gamma\). All new function occurrences come from
the eligible fresh l-template; protected substituted images stay
normal. This replacement is strictly greater than \(u\rho\)
under \(\succ\), so it is not a descent step for the source M.
The lemma proves the chosen local inverse operation, not existence
of a suitable eligible inverse match in every unsolved state.
\end{hblemma}

\begin{hbtheorem}{thm:p12-eb-progress-countercase}{proved=annex:p12.eb.progress-countercase}{A source-scope progress failure with a successful raw inverse-basic branch}
Over \(\mathcal F=\{a,f,g\}\), where a is nullary and f,g
unary, take
\[
 R=\{f(x)\to g(x),\quad f(g(x))\to g(x),\quad g(a)\to a\}.
\]
R respects a ground-total reduction ordering, is left-linear and regular, and
terminates on all open terms. It is ground confluent in this
original signature: every closed term reduces to a. Nevertheless
the open terms \(g(g(y))\) and \(g(y)\) are distinct R-normal
terms and are R-equal. Identity is an R-reduced solution of
\(P=\{g(g(y))=^{?}g(y)\}\), but neither endpoint has a
forward redex. Any R-reduced substitution R-equal to identity
on \(\{y\}\) also maps y to y, so choosing a different normal
image representative does not repair the missing forward step.

In the raw inverse-identity B reading, choose fresh
\(g(v)\approx f(v)\) at the left root. Its matching-first
constraint gives \(v\approx g(y)\), with a normal image.
The retained skeleton becomes \(f(v)=^{?}g(y)\).
Choose the forward rule \(f(g(w))\to g(w)\), fresh w.
First matching Decomposition gives \(g(w)=^{?}g(y)\), then
literal solving gives \(w\approx y\). Constrain on the remaining
\(g(w)=^{?}g(y)\) deletes a literal identity. The terminal
S is \(\{v\approx g(y),w\approx y\}\), so its answer on y
is identity. Normal witness values and basic addresses are preserved
throughout. The first inverse step strictly increases the source
multiset M: \(f(g(y))\succ g(g(y))\). Thus the source's forward-
descent cases do not by themselves exhaust this actual term scope,
even though the raw inverse-basic branch succeeds.
\end{hbtheorem}

\begin{proofboundary}{def:p12-eb-progress-boundary}{What remains after initial normalization and the local history proof}
Initial finite-image normalization is supplied above under the
source's actual reduction-ordering hypothesis. It is not the
unresolved dependency. At the literal fixed-signature ground-
convergence scope, what remains is a global progress argument for
R-equal but different open normal endpoints: for example, a finite
equational-conversion-to-basic construction with a justified measure
and occurrence transport for inverse expansions. The source's
grounding step does not provide that argument without establishing
the requisite convergence in the grounding vocabulary.

The countercase refutes the exhaustion/descent mechanism at that
scope, not raw inverse-basic completeness itself. Its successful
inverse trace prevents transferring a directed failure to the
larger family. Strengthening confluence to the whole open domain
would make a different theorem; a ground-target-only result would
also be a different statement. Neither is silently substituted for
the original source claim. Existing stronger-domain completeness,
terminal normalization and arbitrary-theory hereditary refinements
retain their own proof boundaries.
\cite[Theorem4.13 proof, pp.~497--498]{BaaderSnyder2001}
\end{proofboundary}


\subsection{Strategies and refinements of basic narrowing}

\begin{hbdefinition}{def:p12-eu-narrowing-refinements}{Source strategy refinements and their proof standing}
The source's first four constraint rules form its syntactic subsystem
\(\mathcal S\); composite Solve and Narrow combine one transfer or
LP step with syntactic constraint solving. Their actual helper-MGU,
source-order accumulator and factor conditions are established in
\cref{thm:p12-en-mgu-blocks,thm:p12-en-actual-lift}, including
mandatory matching-first Decomposition and the literal old-constraint
context. Eager A means no rule
outside A is applied while an A rule is enabled anywhere. A
simplification must preserve reduced witnesses on old variables and
strictly decrease the completeness induction measure. The source
additionally supplies Reduce, redex-ordered Narrow with abstraction
of earlier positions, Propagation with measure
\((M,i,n_1,n_2,n_3)\) where i counts nonvariable skeleton symbols,
and Blocking plus the lower-index rule restriction on Narrow.
These refinements require separate checks. The literal source
Propagation rule, and the scope of fresh bindings retained by Reduce,
are distinguished from the corrected combinations below; a blocked
nonreduced branch may still have equational solutions.
\cite[Definitions 4.15--4.16, Theorem 4.17, Proposition 4.18,
\S4.4.3--4.4.4, pp.~499--503]{BaaderSnyder2001}
\end{hbdefinition}

\begin{hbdefinition}{def:p12-eu-eager-strategy}{Reduced-witness simplifications and eager control}
Fix the ordinary convergent R and parameter domain of
\cref{def:p12-eu-narrowing-standing}. A source simplification carries
the printed restriction condition, for old state \(P;C;S\) and
new state \(P';C';S'\):
\[
 \theta\text{ is an R-reduced solution of }P';C';S'
 \quad\Longleftrightarrow\quad
 \theta|_{\operatorname{Var}(P,C,S)}
       \text{ is an R-reduced solution of }P;C;S.
\]
It also requires strict descent relative to those witnesses. For
fresh-variable operations this printed biconditional must be qualified
by the specified extension, as explained below. The checked formulation
carries
solutions bijectively on old variables, with any fresh variables given
their specified literal extensions, preserves reduced witnesses, and
decreases its completeness measure. ``Eager A'' means that a rule
outside A is forbidden while any A rule applies anywhere.
LP's mandatory matching-first Decomposition remains in force during
all eager choices. For fresh-variable rules, extensions must satisfy
their new literal binding; an arbitrary assignment to a fresh variable
is not a witness. The source's Definition 4.15 restriction biconditional
requires this qualification when its Propagation argument introduces
a fresh variable and supplies its particular image.
For B use \((M,i,n_1,n_2,n_3)\): M is the multiset of instantiated
skeleton terms, i the number of nonvariable skeleton symbols, and the
last three are the syntactic constraint counters. For N use
\((M,n_1,n_2,n_3,i)\), an editorial order retaining the same
objects: N Elimination can increase skeleton size while strictly
decreasing its constraint-variable counter.

Corrected Propagation requires u itself to be a nonvariable skeleton
term, as well as the source condition \(u\sigma_S=s\) nonvariable
with s occurring in a solved right side. Reduce matches a fresh
\(l\to r\) to \(u\sigma_S\) by assigning only fresh rule
variables: \(l\gamma=u\sigma_S\). The literal source variant
retains \([\gamma]\) in S; a separate cleaned variant replaces
u by \(r\gamma\) and discards those now-isolated fresh match
bindings. They occur in no old data, no original input variable,
and no fully substituted replacement. This projection is explicit,
not a silent alteration of the source's state.
\cite[Definitions 4.15--4.16 and Reduce, pp.~500--501]{BaaderSnyder2001}
\end{hbdefinition}

\begin{hblemma}{thm:p12-eu-local-simplifications}{proved=annex:p12.n-c-035.5.local-simplifications}{The actual corrected simplification checks}
Constraint simplification, constructor Problem Decomposition, corrected
nonvariable-skeleton Propagation, and cleaned Reduce preserve reduced
witnesses on old variables and strictly decrease the respective measures
in \cref{def:p12-eu-eager-strategy}. Problem Decomposition requires
that its root is absent from every R-left-side root. Propagation's
new image is a normal subterm of an old normal solved image. Cleaned
Reduce constrains no old variable and introduces no retained fresh
binding; its ground instance is an actual decreasing R step.
\end{hblemma}

\begin{hbtheorem}{thm:p12-eu-eager-corrected}{proved=annex:p12.n-c-035.5.eager-corrected}{Eager corrected simplification preserves complete B/N search}
Adding the checked simplifications to B or N and performing them
eagerly preserves coverage of every original E-unifier. In particular,
eager cleaned Reduce together with eager corrected Propagation is
complete under the declared convergence hypothesis. This is an
explicit corrected combination; it is not an unqualified completeness
assertion for the unmodified literal pair of source rules.
The source's abstract eager-simplification result applies whenever
its actual local reduced-witness and descent conditions are met.
\cite[Theorem 4.17, p.~500]{BaaderSnyder2001}
\end{hbtheorem}

\begin{hbtheorem}{thm:p12-eu-eager-reduce-literal}{proved=annex:p12.n-c-035.5.eager-reduce-literal}{Literal Reduce alone can retain harmless isolated bindings}
Without using its isolated fresh match bindings as reduced-term
certificates or rejecting them by Blocking, the literal source Reduce
variant can be performed eagerly in B/N while retaining complete
coverage on the original problem variables. Fresh solved match domains
are absent from all active constraints, skeletons and other right sides;
they may have nonnormal images without constraining the old solution.
Thus they need not satisfy the stronger reduced-witness condition used
for the corrected combined strategy.
\cite[Proposition 4.18, p.~501]{BaaderSnyder2001}
\end{hbtheorem}

\begin{hbexample}{def:p12-eu-propagation-loop}{Why the literal propagation side condition needs repair}
With R empty, start from
\(\{x=^{?}g(b),f(x)=^{?}f(g(b))\}\). After constraining and
eliminating the first equation, S contains \(x\approx g(b)\).
The printed condition permits propagation at the \emph{variable}
occurrence x: \(x\sigma_S=g(b)\) is nonvariable. It creates
\(y\approx g(b)\) and replaces x by y; the same step creates z
from y, and so on. After any nonvariable occurrences are also
abstracted, these aliases remain enabled. An eager literal propagation
strategy can never hand that nonempty problem back to solving; its
claimed i coordinate does not decrease at those variable replacements.
The nonvariable-skeleton condition prevents the cycle while preserving
the intended removal of actual constructor occurrences.
\cite[Propagation and its measure, p.~502]{BaaderSnyder2001}
\end{hbexample}

\begin{hbexample}{def:p12-eu-reduce-certificate-failure}{A fresh match image is not automatically a reduced certificate}
Let \(R=\{f(z)\to a,g(b)\to c\}\) and
\(P=\{f(g(y))=^{?}a,h(g(y))=^{?}h(c)\}\).
The original solution \(y\mapsto b\) is reduced. Literal Reduce
at the first equation retains the fresh binding \(z\approx g(y)\).
Its image under that witness is the reducible \(g(b)\), so this
binding cannot certify that another \(g(y)\) occurrence is reduced.
Eager literal Propagation can turn the second equation into
\(h(w)=^{?}h(c)\), with \(w\approx g(y)\), hiding its only
basic redex. Alternatively narrowing g(y) first then assigning y=b
makes that retained match binding symbolically reducible and Blocking
rejects it. Cleaned Reduce discards the isolated z-binding; narrowing
the second equation then correctly returns \(y\mapsto b\).
This distinguishes the sound source Reduce step from the stronger
certificate invariant needed by the combined strategies.
\end{hbexample}

\begin{hbdefinition}{def:p12-eu-redex-ordered-strategy}{Finite redex-ordered abstraction and lower-index selection}
Order the positions of the current equation by a strict order
\(\prec_{\rm red}\) extending proper-subterm priority: descendants
precede ancestors. For a completeness witness select a minimal
ground redex, and the lowest indexed rule matching that redex.
Redex-ordered Narrow performs the usual LP/constraint-solving block,
then abstracts the earlier retained \emph{context} positions, outside
the replaced hole. Removed positions inside the old redex are not
positions of the new context. Visit each prescribed original position
at most once; an already-variable occurrence needs no further abstraction.
Equivalent implementation abstracts maximal retained earlier regions,
which removes their descendant positions as well. This finite pass is
not an unrestricted eager loop creating fresh aliases repeatedly.

The final lower-index test is literal matching:
\(l\sigma_S\nu\) must not be an instance of any left side of
lower rule index. It is not an E-unifiability test. Blocking checks
reducibility of solved terms or proper subterms of constraints; the
proof branch keeps precisely the normality certificates that justify
those checks, including at the intermediate matching-constraint stage.
\cite[\S4.4.3--4.4.4, pp.~501--503]{BaaderSnyder2001}
\end{hbdefinition}

\begin{hbtheorem}{thm:p12-eu-ordered-narrow-blocking}{proved=annex:p12.n-c-035.5.ordered-narrow-blocking}{Redex order, finite abstraction, Blocking and lower-index completeness}
Under the declared convergence hypothesis, replacing Narrow by the
redex-ordered block with the finite abstraction convention preserves
complete original-variable coverage. Corrected eager simplifications
may be included. Eager Blocking and the source lower-index restriction
do not remove the constructed complete witness branch: its solved
images are normal, its constraint proper subterms are normal, and its
selected rule is the first indexed rule matching the ground redex.
This proves the actual side-condition and lifting cases, not merely
the assertion that these strategies are ``consistent.''
\end{hbtheorem}

\subsection{ACU, ACUI, and Abelian groups}

\begin{hbdefinition}{def:p12-eu-three-theories}{Three exact source theories and coordinate domains}
For binary \(f\), constant \(e\), and unary \(i\) in AG, put
\[
\begin{aligned}
 ACU&=\{f(x,y)\approx f(y,x),\ f(f(x,y),z)\approx f(x,f(y,z)),\
                          f(x,e)\approx x\},\\
 ACUI&=ACU\cup\{f(x,x)\approx x\},\\
 AG&=ACU\cup\{f(x,i(x))\approx e\}.
\end{aligned}
\]
An ACU term has a multiplicity vector over \(\mathbb N\); an ACUI
term has a support vector over \(BS=\{0,1\}\), with \(\lor\)
addition and \(\land\) multiplication; an AG term has a signed
multiplicity vector over \(\mathbb Z\). Zero is represented by
\(e\), and negative coordinates in AG by inverses. Write
\(\operatorname{coef}_q(t)\) for the coordinate at a generator \(q\).
The source's ACU notation \(|t|_q\) is an occurrence count local to
this presentation, not cardinality of a term.
\cite[\S5.1, pp.~504--506]{BaaderSnyder2001}
\end{hbdefinition}

\begin{hblemma}{thm:p12-eu-coordinate-models}{proved=annex:p12.n-c-035.5.coordinate-models}{Coordinates characterize theory equality}
For elementary terms and terms with free constants added, equality modulo
ACU, ACUI, or AG is exactly equality of the respective generator
coordinates. Every finite-support vector is represented by a term.
Substitution composition acts by matrix multiplication over the respective
coefficient domain \(\mathbb N\), \(BS\), or \(\mathbb Z\).
\end{hblemma}

\begin{hbmethod}{def:p12-eu-homogeneous-construction}{The elementary finite-generator construction}
List input variables \(x_1,\ldots,x_n\) and equations
\(s_j=^{?}_E t_j\), \(1\leq j\leq m\). Let \(M_1,M_2\)
be the \(m\times n\) matrices of left and right variable coordinates.
For the appropriate coefficient semiring \(S\), put
\[
                U(M_1,M_2)=\{v\in S^n:M_1v=M_2v\}.
\]
Choose finite generators \(v_1,\ldots,v_k\), so every solution is
\(\sum_jv_js_j\), with \(s_j\in S\). Introduce distinct fresh
\(z_1,\ldots,z_k\) and make the \(z_j\)-coordinate of
\(x_i\sigma\) equal to \((v_j)_i\). Zero rows give \(e\);
negative coefficients in AG give inverse terms. The substitution matrix
has columns \(v_j\). Finding generators and constructing \(\sigma\)
are distinct steps; existence of generators does not provide a stopping
criterion for naive unbounded basis enumeration.
\cite[pp.~505--506 and Theorem 5.4, p.~511]{BaaderSnyder2001}
\end{hbmethod}

\begin{hbtheorem}{thm:p12-eu-homogeneous-completeness}{proved=annex:p12.n-c-035.5.homogeneous-completeness}{Generators give a most general elementary answer}
The preceding substitution is a most general \(E\)-unifier. For ACU,
the componentwise minimal nonzero nonnegative solutions form a finite
generating set. For ACUI all nonzero solution vectors form a finite
generating set. For AG a supplied finite integer generating set suffices;
factors use integer, not just nonnegative, coefficients. Integer kernels
also have finite generators. These assertions establish unitary elementary
unification in all three theories, including
the empty-basis answer sending all input variables to \(e\).
They assert neither termination of unrestricted narrowing nor termination
of a naive search for the ACU generators.
\end{hbtheorem}

\begin{hbexample}{def:p12-eu-source-elementary-example}{Three versions of the source's equation}
For \(f(x,f(x,y))=^{?}_E f(z,f(z,z))\), the ACU equation is
\(2x+y=3z\) in \(\mathbb N\), with minimal nonzero solutions
\[
 v_1=(1,1,1),\qquad v_2=(0,3,1),\qquad v_3=(3,0,2).
\]
The three corresponding unifiers are
\[
\begin{aligned}
 \sigma_1&=\{x\mapsto z_1,y\mapsto z_1,z\mapsto z_1\},\\
 \sigma_2&=\{x\mapsto e,y\mapsto f(z_2,f(z_2,z_2)),z\mapsto z_2\},\\
 \sigma_3&=\{x\mapsto f(z_3,f(z_3,z_3)),y\mapsto e,z\mapsto f(z_3,z_3)\}.
\end{aligned}
\]
None alone is most general. Their coordinatewise combination is
\[
\begin{aligned}
 x\sigma&=f(z_1,f(z_3,f(z_3,z_3))),\\
 y\sigma&=f(z_1,f(z_2,f(z_2,z_2))),\\
 z\sigma&=f(z_1,f(z_2,f(z_3,z_3))).
\end{aligned}
\]
Sending the other two fresh variables to \(e\) recovers each
\(\sigma_j\). Every nonnegative solution is a sum of the three vectors:
subtract \(v_1\) while both first coordinates are positive; with one
zero use \(v_2\) or \(v_3\). The source's further example has
\((1,4,2)=v_1+v_2\) at \(z'\) and \((3,0,2)=v_3\) at
\(z''\); it factors through \(\sigma\) by
\(z_1\mapsto z',z_2\mapsto z',z_3\mapsto z''\).

For ACUI the equation becomes \(x\lor y=z\), generated over
\(BS\) by \((1,0,1)\) and \((0,1,1)\). Its most general answer is
\(\sigma'=\{x\mapsto z_1,y\mapsto z_2,z\mapsto f(z_1,z_2)\}\).
Assigning \(z_2\mapsto z_1\) gives \(\sigma_1\) modulo ACUI,
using idempotence, not ACU equality.
For AG the equation is \(2x+y=3z\) over \(\mathbb Z\), generated
by \((0,3,1)\) and \((1,-2,0)\), since every solution is
\(z(0,3,1)+x(1,-2,0)\). The answer is
\[
 \sigma''=\{x\mapsto z_2,
 y\mapsto f(z_1,f(z_1,f(z_1,f(i(z_2),i(z_2))))),z\mapsto z_1\}.
\]
The minus sign is essential: \((1,2,0)\) does not solve this equation.
\cite[\S5.1, pp.~504--506]{BaaderSnyder2001}
\end{hbexample}

\begin{hbmethod}{def:p12-eu-constant-construction}{Free constants give inhomogeneous systems}
Let \(C\) be the input free constants. For each \(c\in C\),
its input left/right coordinate vectors \(u_{1,c},u_{2,c}\) impose
\[
                    M_1w_c+u_{1,c}=M_2w_c+u_{2,c}
\]
on its coordinate vector in variable images. Choose a finite particular
set \(P_c\) such that every inhomogeneous solution is \(p+v\),
\(p\in P_c\), \(v\in U(M_1,M_2)\). For every choice of one
\(p_c\) per constant, add those fixed constant coordinates to the
homogeneous generating substitution. Collect all such answers. If
one \(P_c\) is empty there is no solution. Constants are fixed
throughout, not substituted as variables.
\cite[\S5.1, pp.~507--508]{BaaderSnyder2001}
\end{hbmethod}

\begin{hbtheorem}{thm:p12-eu-constant-completeness}{proved=annex:p12.n-c-035.5.constant-completeness}{The inhomogeneous construction is finite complete}
Under the generating and particular-set hypotheses just stated, the
constructed set is finite complete. For ACU the minimal inhomogeneous
solutions give \(P_c\); for ACUI a finite choice exists since
\(BS^n\) is finite; for AG any one particular solution suffices
for each nonempty inhomogeneous set, using integer homogeneous factors.
Thus AG with constants is unitary when solvable. ACU and ACUI with
constants can require several incomparable answers. Arbitrary complete
choices of basis and particular solutions do not themselves ensure minimality.
\end{hbtheorem}

\begin{hbexample}{def:p12-eu-source-constant-example}{Two incomparable ACU answers}
For \(f(x,f(x,y))=^{?}_{ACU}f(a,f(z,f(z,z)))\), with \(a\) free,
the homogeneous equation is \(2x+y=3z\), and the \(a\)-coordinate
equation is \(2x+y=3z+1\). Its minimal solutions are \((0,1,0)\)
and \((2,0,1)\). With the explicit terms \(x\sigma,y\sigma,z\sigma\)
of the elementary example, the source's two answers are
\[
\begin{aligned}
 \alpha&=\{x\mapsto x\sigma,y\mapsto f(a,y\sigma),z\mapsto z\sigma\},\\
 \beta&=\{x\mapsto f(a,f(a,x\sigma)),y\mapsto y\sigma,
                              z\mapsto f(a,z\sigma)\}.
\end{aligned}
\]
These names denote the already displayed terms, not recursive image
substitution. Every instance of \(\alpha\) has an \(a\) in its
\(y\)-image, while \(\beta\) with its fresh variables sent to \(e\)
does not. Every instance of \(\beta\) has two \(a\)'s in \(x\),
while that specialization of \(\alpha\) has none. Neither covers the
other. Completeness follows from the construction theorem, so they
are minimal complete. The source's three-constant count \(2\cdot3\cdot5=30\)
presupposes independent respective minimal particular choices of sizes
\(2,3,5\), not merely three constants.
\cite[pp.~507--508]{BaaderSnyder2001}
\end{hbexample}

\begin{proofboundary}{def:p12-eu-ac-not-acu}{AC is not ACU}
AC has no unit axiom or empty-product image. Switching off generators
by assigning \(e\) is unavailable. The source records AC as finitary
in all three regimes; ACU as unitary elementarily and finitary with
constants and generally; and AG as unitary elementarily and with
constants but finitary generally. The elementary matrix construction
does not by itself account for arbitrary free-symbol subterms or for
the nonempty-image conditions of AC. The pure-AC construction in
\cref{sec:p12-equational-ac} supplies the nonempty-image algorithm,
with finite complete and minimal families. Its AC-LCR bridge and
the actual disjoint-theory algorithm in
\cref{thm:p12-comb-ac-lcr,sec:p12-equational-combination} extend
that independent route to arbitrary finite nested/multiple AC symbols
with free symbols and constants; they do not identify AC with ACU.
\cite[\S3.4, pp.~484--486]{BaaderSnyder2001}
\end{proofboundary}

\subsection{Pure AC: Positive Images and Finite Complete Families}
\label{sec:p12-equational-ac}

The source distinguishes AC from ACU and refers to the AC algorithms
of Stickel and others; its semantic \S5 develops the monoidal examples
ACU, ACUI and AG, not an implicit unit for AC. The variable-only AC
subset construction below is explicitly present in the separately
identified 1975 paper. Its extension here to simultaneous equations and
fixed free constants has full algebraic proofs and a finite basis algorithm.
It does not identify that 1975 manifestation with the 1981 journal paper.
\cite[\S3.4 and \S5.1, pp.~484--486,504--508]{BaaderSnyder2001}
\cite[pp.~72--75]{Stickel1975AC}

\begin{hbdefinition}{def:p12-ac-input}{The unit-free term domain and input matrices}
Let \(AC=\{f(f(x,y),z)\approx f(x,f(y,z)),f(x,y)\approx f(y,x)\}\),
with binary f and no designated unit or unit axiom. A flat AC problem is a finite set of
equations whose terms use f, variables, and free constants only; nested
f nodes are flattened. A single variable or constant is a legal term.
An empty product is not. In this construction no further positive-arity
symbol is in the term signature. In examples, multiargument f notation
abbreviates a binary fold, not a new arity; a one-atom product is that
atom itself. All substitution images are finite nonempty terms, and
constants are fixed. Compare answers only on the original
\(X=\operatorname{Var}(\Gamma)\), using the already declared AC-instance
preorder. Enumerate \(X=(x_1,\ldots,x_n)\), and the input constant
set C. Let \(M_L,M_R\) be the variable-multiplicity matrices and
\(u_{L,c},u_{R,c}\) each constant's multiplicity columns. Put
\(D=M_L-M_R\) and \(d_c=u_{R,c}-u_{L,c}\).
The homogeneous columns solve \(Dv=0,v\in\mathbb N^n\);
the c-columns solve \(Dw=d_c,w\in\mathbb N^n\).
Zero is an internal multiplicity vector, never an AC term image.
\end{hbdefinition}

\begin{hblemma}{thm:p12-ac-multisets}{proved=annex:p12.ac.multisets}{The exact unit-free coordinate model and cancellation}
For terms over f, variables and free constants, AC equality is equality
of their nonempty generator multisets. Substitution composition acts on
counts by nonnegative integer matrix multiplication. Common flattened
arguments can be cancelled without losing solutions. An empty residual
argument list after cancellation is only an internal balance constraint;
it neither introduces a unit nor changes the nonempty-image condition.
\cite[pp.~72,74]{Stickel1975AC}
\end{hblemma}

\begin{hbmethod}{def:p12-ac-single-bound}{The source's finite single-equation basis search}
After cancelling common variable counts, a single homogeneous equation
has positive coefficients on its remaining two sides:
\(\sum_{i=1}^{m}a_ix_i=\sum_{j=1}^{n}b_jy_j=k\).
If both sides are present, let
\(L=\max_{i,j}\operatorname{lcm}(a_i,b_j)\) and
\(K=\max(m,n)L\). Enumerate all nonzero nonnegative solutions
with balance k at most K; each coordinate is bounded by K divided
by its coefficient. Retain only the componentwise minimal solutions.
Variables cancelled completely are unconstrained, not assigned zero.
A remaining one-sided positive equation has only the zero solution.
\cite[Lemma and algorithm, pp.~73--75]{Stickel1975AC}
\end{hbmethod}

\begin{hblemma}{thm:p12-ac-single-bound}{proved=annex:p12.ac.single-bound}{The actual lcm bound proves finite generation}
Every nonnegative solution of the positive-coefficient equation in
\cref{def:p12-ac-single-bound} is an additive combination of solutions
with k at most K. Every minimal nonzero solution is therefore within
that finite search. The retained solutions generate all solutions.
\end{hblemma}

\begin{hbmethod}{def:p12-ac-system-basis}{A finite simultaneous-equation basis algorithm}
For any integer matrix A with n columns, enumerate every nonempty column
support S. Compute \(\ker A_S\) by rational Gaussian elimination.
If its dimension is one and a spanning vector has all coordinates
of one strict sign, choose its positive primitive integer vector,
extended by zeros off S. Collect the distinct vectors \(r_1,\ldots,r_t\).
Put \(B_i=\sum_j(r_j)_i\), taking B zero if no vector exists.
Enumerate the finite integer box \(0\leq v_i\leq B_i\), retaining
its nonzero solutions of \(Av=0\). Remove any vector having a
distinct nonzero retained solution componentwise below it; call the
remaining set \(H(A)\). All operations use exact rational/integer
arithmetic. This is an editorial system-level implementation, not a
claim that this ray/box procedure is Stickel's printed single-equation
search or the source's cited incremental algorithms.
\end{hbmethod}

\begin{hbtheorem}{thm:p12-ac-system-basis}{proved=annex:p12.ac.system-basis}{The finite search returns exactly the minimal homogeneous generators}
The preceding algorithm terminates and returns precisely the minimal
nonzero nonnegative integer solutions of \(Av=0\). They generate
all nonnegative integer solutions by addition. The ray set alone may
not generate the integer solutions; the bounded integer remainder
search is essential. No unbounded enumeration with an unstated stopping
criterion is used.
\end{hbtheorem}

\begin{hbmethod}{def:p12-ac-particulars}{Fixed constants through an augmented homogeneous system}
For \(Dw=d_c\), apply the basis algorithm to
\(\widetilde D_c=[D\mid -d_c]\). Select the basis vectors whose
last coordinate is one and project their first n coordinates; this
finite set is \(P_c\). Compute \(H=H(D)=\{h_1,\ldots,h_k\}\)
once. A missing \(P_c\) signals that no c-column is possible.
\end{hbmethod}

\begin{hblemma}{thm:p12-ac-particulars}{proved=annex:p12.ac.particulars}{All and only minimal affine columns are found}
The projected set \(P_c\) is exactly the componentwise minimal
nonnegative solutions of \(Dw=d_c\). Every solution is
\(p+v\), with \(p\in P_c\) and v a homogeneous solution.
This includes \(d_c=0\), where \(P_c=\{0\}\); constants
cancelled from the equations can still occur in later parameter images.
\end{hblemma}

\begin{hbmethod}{def:p12-ac-extraction}{Positive-support branch construction and canonical extraction}
Choose one \(p_c\in P_c\) for each input constant. Enumerate
every subset \(I\subseteq\{1,\ldots,k\}\). Retain it only if
\[
             \sum_{c\in C}(p_c)_i+\sum_{j\in I}(h_j)_i>0
                    \qquad(1\leq i\leq n).
\]
Use a distinct fresh variable \(z_j\) for each selected j. Its
image in \(x_i\sigma_{p,I}\) has \((h_j)_i\) occurrences,
and c has \((p_c)_i\) occurrences. Assemble that nonempty multiset
as one atom when its size is one and as a fixed right-associated
binary f term otherwise. All variables outside X are fixed.
The empty I is permitted when the constants alone make every image
nonempty. A zero row is rejected, never represented by an invented
unit or \(f()\). An input constant named e is an ordinary atom,
not a zero-multiplicity image.
Canonical integer vectors, fixed original-variable order, sorted
constants and fixed fresh-variable choices deduplicate literal outputs.
For a supplied target unifier, its factor can also be computed:
choose the first \(p_c\leq w_c\) in each finite \(P_c\), then
repeatedly subtract the first \(h_j\leq v\) from each homogeneous
residual column v until it is zero. Record the subtraction counts as
parameter-image multiplicities. Each subtraction strictly lowers the
coordinate sum, so this factor construction terminates as well.
\end{hbmethod}

\begin{hbtheorem}{thm:p12-ac-coverage}{proved=annex:p12.ac.coverage}{The complete finite pure-AC algorithm}
For a flat finite AC problem, the basis, augmented-column and positive-
support procedures terminate and return a finite complete set of AC
unifiers. Every arbitrary AC unifier is covered by a returned answer
whose factor assigns a nonempty term to each selected fresh variable.
Zero coefficients decide which generators are omitted, not which
parameters are assigned empty terms. Empty input returns identity.
For n=0 compare the ground input multisets directly, returning identity
if they agree and the empty answer set otherwise.
\end{hbtheorem}

\begin{hbtheorem}{thm:p12-ac-minimality}{proved=annex:p12.ac.minimality}{Minimality and deduplication are established, not assumed}
With the actual minimal H and minimal \(P_c\) computed above,
distinct admissible pairs \((p,I)\) give incomparable answers
on X, up to the fixed canonical fresh naming. Thus the deduplicated
output is a minimal complete set, not merely complete. A nonminimal
generating set would not justify this conclusion. No uniqueness of
factor substitutions is asserted.
\end{hbtheorem}

\begin{hbexample}{def:p12-ac-five-answers}{The source's ACU equation has five unit-free answers}
For \(f(x,f(x,y))=^{?}_{AC}f(z,f(z,z))\), H consists of
\(h_1=(1,1,1),h_2=(0,3,1),h_3=(3,0,2)\).
Exactly five subsets cover every variable row:
\(\{1\},\{2,3\},\{1,2\},\{1,3\},\{1,2,3\}\).
For example \(\{1\}\) yields x=y=z=\(z_1\), whereas
\(\{2,3\}\) yields \(x=f(z_3,f(z_3,z_3))\),
\(y=f(z_2,f(z_2,z_2))\), \(z=f(z_2,f(z_3,z_3))\).
Singleton 2 is rejected for its empty x-image; singleton 3 for its
empty y-image. The five answers are minimal complete by the theorem.
ACU instead allows switching off parameters by e and combines them
into one answer. ACUI changes counts to supports, and AG allows signed
counts; neither operation belongs to this AC algorithm.
The two positive cone rays are h2 and h3; h1 is an indispensable
integer remainder, not an integer sum of those rays.
\cite[pp.~505--506]{BaaderSnyder2001}
\end{hbexample}

\begin{hbexample}{def:p12-ac-stickel-69}{The actual 1975 seven-generator and constant example}
The source variable equation is
\(2x_1+x_2+x_3=2y_1+y_2\). In source order its seven minimal
vectors are
\[
\begin{gathered}
(0,0,2,1,0),\ (0,1,0,0,1),\ (0,0,1,0,1),\ (0,1,1,1,0),\\
(0,2,0,1,0),\ (1,0,0,0,2),\ (1,0,0,1,0).
\end{gathered}
\]
Exactly 69 of its 128 subsets have positive sums in all five rows.
For \(f(x,x,y,a)=^{?}_{AC}f(b,b,z)\), constant a is the
third row and b the fourth. A constant row must contain exactly
one parameter occurrence, not an f-product of several occurrences.
This leaves six subsets; any using vector 4 forces the same parameter
to equal distinct a and b, leaving four answers. The direct fixed-
constant procedure computes the same minimal family, with fresh p,q:
\[
\begin{aligned}
 &(x,y,z)=(b,p,f(a,p)),\\
 &(x,y,z)=(f(b,q),p,f(a,f(p,f(q,q)))),\\
 &(x,y,z)=(q,f(b,b),f(a,f(q,q))),\\
 &(x,y,z)=(q,f(b,f(b,p)),f(a,f(p,f(q,q)))).
\end{aligned}
\]
Here \(H(D)=\{(0,1,1),(1,0,2)\}\),
\(P_a=\{(0,0,1)\}\), and
\(P_b=\{(1,0,0),(0,2,0)\}\). These yield the four positive
branches without treating a or b as substitutable variables.
\cite[pp.~73--75]{Stickel1975AC}
\end{hbexample}

\begin{hbexample}{def:p12-ac-degeneracies}{Cancellation, constants and impossible positive images}
For \(x=^{?}_{AC}f(x,x)\), the homogeneous kernel is zero;
its only numerical column would make the image empty, so the answer
set is empty. ACU has \(x\mapsto e\); idempotent ACUI has identity.
For \(f(x,a)=^{?}_{AC}f(x,b)\) with a,b distinct, cancellation
leaves incompatible fixed constant columns, so there is no answer.
For \(f(x,y)=^{?}_{AC}f(a,b)\), the two answers are
\((x,y)=(a,b)\) and \((b,a)\). The other numerical allocations
leave one image empty and are rejected. ACU additionally permits
\((f(a,b),e)\) and \((e,f(a,b))\).
An identically zero variable column is unconstrained and has a unit
coordinate generator; it must still receive a nonempty variable image.
For \(f(x,x)=^{?}_{AC}a\), the augmented kernel has generator
\((1,2)\), but no last-coordinate-one generator; there is no
integer particular column, even though a rational half-column exists.
Duplicating an equation does not change the kernel or answer classes.
All these cases are handled by the stated finite procedure.
\end{hbexample}

\begin{proofboundary}{def:p12-ac-general-boundary}{Flat construction and general-input source boundaries}
This procedure covers every flat f/variable/constant system. It does
not silently turn a free-function-rooted input argument into a constant.
Stickel 1975 generalizes such arguments to fresh variables and then
reconciles their values by recursive unification, imposing single-atom
constraints where required. Its general-term proof invokes recursive
calls of lesser complexity. The paper expressly states termination,
soundness and completeness for that general-term algorithm; the flat
proof is not presented as establishing that stronger source statement.
The historical recursive algorithm's unrestricted nested-input
termination argument and the distinct 1981 journal presentation are
not supplied by the flat matrix proof here. An independent full
route for arbitrary finite nested/multiple AC symbols and free
symbols/constants is supplied by the AC-LCR bridge and actual ordered
combination construction in
\cref{thm:p12-comb-ac-lcr,sec:p12-equational-combination}, with its
own complete proofs and terminating component procedures. That route
does not discharge these distinct historical-source presentations.
The source's \S5 similarly states that its semantic approach does
not directly solve general free-function inputs; the later combination
treatment supplies a justified bridge, not an unmentioned extension
of that semantic method.
\cite[pp.~73--76]{Stickel1975AC}
\cite[pp.~504,512--513]{BaaderSnyder2001}
\end{proofboundary}


\subsection{The semiring bridge}

\begin{hbdefinition}{def:p12-eu-monoidal-commutative}{Monoidal syntax and commutative free-algebra structure}
A monoidal theory has one binary \(f\), one constant \(e\), and
otherwise unary symbols; \(f\) is associative and commutative with
unit \(e\), and each unary \(h\) preserves \(f\) and \(e\).
A commutative theory in the source's broader sense has semi-additive
\(C_{\operatorname{Sig}(E)}(E)\): a zero object and finite coproducts
that are also products with their biproduct maps. In algebraic terms the
initial free algebra has one element, and the free algebra on disjoint
finite generator sets has the canonical direct-product decomposition.
``Commutative'' is this categorical property, not commutativity of every
primitive operation in an arbitrary signature presentation.
\cite[Definitions 5.1--5.2, pp.~508--509]{BaaderSnyder2001}
\end{hbdefinition}

\begin{hbdefinition}{def:p12-eu-semiring}{One-generator endomorphisms}
For a commutative theory let \(S_E\) be the endomorphisms of its free
one-generator algebra. Multiplication is source-order composition;
addition is biproduct addition, equivalently pointwise application of
the implicit monoid operation \(*\):
\((\sigma+\tau)(x)=\sigma(x)*\tau(x)\). Zero is the zero map,
and one is identity. Multiplication need not be commutative.
The matrix of an arrow has one row per source generator and one column
per target generator; \(M_{\sigma\tau}=M_\sigma M_\tau\).
For ACUI, \(x_1\mapsto f(y_1,y_2),x_2\mapsto y_2\) gives
\(\left(\begin{smallmatrix}1&1\\0&1\end{smallmatrix}\right)\).
The three example semirings are \(\mathbb N\), \(BS\), and
\(\mathbb Z\), respectively. \cite[\S5.3, p.~510]{BaaderSnyder2001}
\end{hbdefinition}

\begin{hbtheorem}{thm:p12-eu-semiring-criterion}{proved=annex:p12.n-c-035.5.semiring-criterion}{The finite-generation criterion with its matrix hypothesis}
For a nontrivial commutative theory with this matrix representation,
elementary unification is unitary exactly when every
\(U(M_1,M_2)=\{v\in S_E^n:M_1v=M_2v\}\), for parallel
\(m\times n\) matrices, is finitely generated as a right semimodule.
Its generating columns give a most general answer.
If elementary unification is unitary, unification with constants is
at most finitary exactly when every inhomogeneous solution set
\(\{w:M_1w+u_1=M_2w+u_2\}\) is a finite union of translates
\(p+U(M_1,M_2)\). Constant-respecting matrices have block form
\(\left(\begin{smallmatrix}M^h&M^i\\0&I\end{smallmatrix}\right)\).
The conditions on a solution \(\delta\) are
\[
 M_\sigma^hM_\delta^h=M_\tau^hM_\delta^h,
 \qquad M_\sigma^hM_\delta^i+M_\sigma^i
       =M_\tau^hM_\delta^i+M_\tau^i.
\]
\cite[Theorems 5.4--5.5, pp.~511--512]{BaaderSnyder2001}
\end{hbtheorem}

\subsection{Disjoint-Theory Combination with Linear Constant Restrictions}
\label{sec:p12-equational-combination}

The combination method in source \S6 is the Baader--Schulz method:
purification, identification of shared variables, signature labels, and a
linear order inducing constant restrictions. Its hypotheses are nontrivial
component theories over pairwise disjoint signatures and solvers for
\emph{linear-constant-restricted} problems. Neither regularity nor
collapse-freeness is assumed. The source gives the ordered assembly
argument and a normalization-based completeness sketch. The full
completeness proof below instead constructs the combined free algebra;
it is an explicit equational specialization of the alternative algebraic
route identified by the chapter, not a transcription of the source's
ordered-completion proof or of the quasi-free-structure fibering algorithm.
\cite[\S6, pp.~513--519]{BaaderSnyder2001}
\cite[author manuscript \S\S4.2--4.4, pp.~26--35]{BaaderSchulz1998Combination}

\begin{hbwarning}{def:p12-comb-source-context}{The source's disjointness and historical-algorithm boundaries}
The source's applications include distinct AC operations and combinations
of ACI set union with associative list append, for example sets of
lists. Separate unifiers and separate decision procedures are two
different combination tasks. One-symbol algorithms do not themselves
handle those mixed inputs.

Disjointness cannot be removed from the general theorem. In the source's
notation, \(D^l_{f,g}\) is
\(f(x,g(y,z))\approx g(f(x,y),f(x,z))\),
\(D^r_{f,g}\) is
\(f(g(y,z),x)\approx g(f(y,x),f(z,x))\), and
\(D_{f,g}=D^l_{f,g}\cup D^r_{f,g}\).
For problems with constants, both one-sided theories are reported
unitary, whereas their union is reported infinitary; their signatures
overlap. Similarly, the source reports decidability with constants for
\(D_{f,g}\) and for associative \(A_g\), but undecidability with
constants for \(D_{f,g}\cup A_g\). These historical results
are source-reported counterexamples to unqualified nondisjoint
combination, not full distributivity/undecidability proofs in this entry.
Special nondisjoint methods require additional conditions.
\cite[\S3.4 and p.~513]{BaaderSnyder2001}

Earlier source-listed methods imposed collapse-freeness or regularity:
collapse-free excludes an identity between a variable and a nonvariable
term; regular requires the same variable set on both sides. The source
distinguishes Schmidt--Schau\ss's 1989 method, requiring finitary
unification with constants and finitary constant elimination, and
Boudet's 1993 efficiency refinement with the same requirements.
Their actual constant-elimination constructions are separate obligations,
not supplied by naming them. The Baader--Schulz LCR interface combines
the two needs into one restricted local problem and can also combine
decision procedures without enumerating unifiers.
\cite[p.~514]{BaaderSnyder2001}
\end{hbwarning}

\begin{hbdefinition}{def:p12-comb-domain}{Disjoint signatures and the component interface}
Fix \(n\geq1\) theories \(E_1,\ldots,E_n\), with pairwise disjoint
signatures \(\Sigma_1,\ldots,\Sigma_n\), and put \(E=\bigcup_i E_i\).
All operations have finite arity; terms and substitutions are finite.
Nontrivial means that \(E_i\) does not imply the identity \(u\approx v\)
for distinct variables. Thus its free algebra on two or more generators
keeps these generators distinct. Constants belonging to \(\Sigma_i\)
belong only to that signature. Additional free constants, when required,
may be placed in a separate empty-equational-theory signature; constants
shared as interpreted symbols are outside this theorem.

For a component problem with variables \(V_i\), free constants \(K_i\),
and a strict linear order, an LCR answer \(\mu_i\) fixes \(K_i\), solves
every equation modulo \(E_i\), and obeys
\(x<c\Rightarrow c\notin\operatorname{VarConst}(x\mu_i)\).
Here \(\operatorname{VarConst}(t)\) is the set of variable and
designated-free-constant names occurring in t. The occurrence test treats
the designated free constants as atoms,
not as substitutable variables. A complete answer set covers every
permitted solution by \(E_i\)-instantiation on \(V_i\), with the
factor substitutions fixing \(K_i\). This is the local version of
Definition~\ref{def:p12-eu-lcr} and the problem-variable comparison of
Definition~\ref{def:p12-eu-generality}; it is not comparison on all
variables of the infinite variable supply.
\end{hbdefinition}

\begin{hbmethod}{def:p12-comb-purification}{Purification and the reconstruction map}
An equation with different-signature nonvariable roots is replaced by
\(a=^{?}s,\ a=^{?}t\), for fresh a. In each equation with root signature
\(\Sigma_i\), replace a maximal subterm with root in another signature
by a fresh variable b and add \(b=^{?}u\). Repeat on the new defining
equations. Variable--variable equations are eliminated by identifying
their variables, retaining the resulting substitution \(\lambda\) for
reconstruction. Equations with a variable on one side are assigned to
the signature of the other side. A pure problem is obtained:
\(\Gamma=\Gamma_1\cup\cdots\cup\Gamma_n\), where every term in
\(\Gamma_i\) uses \(\Sigma_i\) and variables only.

A polynomial implementation can name every original nonvariable term
node, bottom-up. For a node \(f(t_1,\ldots,t_m)\), add
\(a_t=^{?}f(a_{t_1},\ldots,a_{t_m})\), using the original variable
itself when a child is a variable. Identify the names of the two roots
of each original equation. For a nullary node c use \(a_c=^{?}c\).
This more fully flattened presentation is an editorial implementation
of the same purification, not a different unification theory. It uses
at most one auxiliary per original term node. Store the quotient of
the variable identifications and the original-variable reconstruction
\(\lambda\); no cyclic term substitution is needed to represent it.
If a representative of an original variable disappears from all
remaining equations, retain it as an unconstrained variable by adding
the tautology \(r=^{?}r\) to one chosen component. This harmless
scope retention makes the original-variable reconstruction literal
even for an input consisting only of variable equations.
\end{hbmethod}

\begin{hblemma}{thm:p12-comb-purification}{proved=annex:p12.comb.purification}{Purification preserves all solutions on the original variables}
Every solution of the original problem extends to a solution of its
purified problem. Conversely, every purified solution, composed with
\(\lambda\) and restricted to the original variables, solves the
original problem. The extension/restriction also preserves coverage:
a complete family for the purified problem reconstructs a complete
family for the original problem. Fully flattened purification has
polynomial size and terminates.
\end{hblemma}

\begin{hbmethod}{def:p12-comb-branch}{The finite branch data and local problems}
Let X be the variables occurring in at least two pure blocks. Enumerate
all partitions \(\Pi\) of X, choose one representative for each block,
and let \(\sigma_\Pi\) replace each shared variable by its representative.
Write \(X_\Pi\) for the representatives. Enumerate every label function
\(L:X_\Pi\longrightarrow\{1,\ldots,n\}\) and every strict linear order
\(<\) on \(X_\Pi\). In component i, representatives with label different
from i are \emph{distinct free constants}; the remaining representatives
are variables. Nonshared variables of \(\Gamma_i\) belong to i and
are placed after all shared representatives. Their relative order is
immaterial since there are no nonshared free constants. The induced
restrictions therefore say: a foreign shared constant y may occur
in an own shared variable x's image only if \(y<x\).

The local problem is \(\langle\Gamma_i\sigma_\Pi,L,<\rangle\).
Labels may be assigned to a representative even if it does not occur
in its labelled block: add the tautology \(x=^{?}x\) to that local
problem, so its complete-answer comparison includes x as an
unconstrained input variable. This matters when \(n>2\).
The identity on an unconstrained variable is a solution; freshening its
range gives a fresh parameter. Ignore a branch when any local answer
set is empty. If X is empty there is one empty partition/order/label
choice, not zero branches.
\cite[pp.~515--517, including footnote~21]{BaaderSnyder2001}
\end{hbmethod}

\begin{hbmethod}{def:p12-comb-merge}{Fresh ranges and the actual ordered merge}
For a chosen local answer, rename every substitutable variable appearing
in its range to a fresh parameter. Fix the foreign shared constants.
Choose mutually disjoint parameter sets \(P_i\), also disjoint from all
input variables. This post-renaming is a bijective renaming of variables
in the local term language, so it preserves the equations and LCR tests;
its inverse shows that it does not change local generality. The local
input scope includes the unconstrained owned representatives just
specified; freshening therefore covers their images as well.
Let the resulting answers be \(\mu_i\).

Process shared representatives in increasing order. For x with label i,
start with \(x\mu_i\) and replace every foreign shared constant y in
that image by the already constructed term \(y\sigma\). Keep \(P_i\)
fixed. After all shared variables, do the same for each nonshared
variable using its unique owner's answer. The resulting finite
substitution is denoted
\(\operatorname{Merge}_{L,<}(\mu_1,\ldots,\mu_n)\).
Reconstruct an original answer as
\(\lambda\sigma_\Pi\operatorname{Merge}_{L,<}(\mu_1,\ldots,\mu_n)\),
restricted to the original input variables. The source's notation
\(\sigma_1\cdots\sigma_n\) for a combined answer denotes this
assembly; arbitrary substitution composition is not a replacement
for it. Output a DAG if term sharing is desired; expanded trees remain
finite but can be exponentially larger.
\cite[pp.~517--518]{BaaderSnyder2001}
\end{hbmethod}

\begin{hbtheorem}{thm:p12-comb-soundness}{proved=annex:p12.comb.soundness}{Ordered merge is well defined, terminating and sound}
Every tuple of LCR solutions in a branch has a finite ordered merge.
The merge solves \(\Gamma\sigma_\Pi\) modulo E, and the reconstructed
answer solves the original problem. This assertion requires neither
acyclicity of arbitrary target substitutions nor normalization of their
written images: the finite branch's LCR conditions establish the actual
merge dependency order.
\end{hbtheorem}

\begin{hbmethod}{def:p12-comb-source-witness}{The original normalization-based completeness sketch}
The source's \S6.2 proposes the following branch from a normalized
combined solution \(\tau\): identify x,y exactly when
\(x\tau=_Ey\tau\); assign the signature label of the root of
\(x\tau\), with label 1 for a variable image; and extend the strict
subterm relation between distinct image terms to a linear order.
For component i, replace alien subterms by new parameters consistently,
using the same parameter for E-equivalent aliens. With an arbitrary
unnormalized solution this abstraction need not solve the pure block,
so these three choices alone are not a completeness proof.

The original proposed normalization is by the ordered system obtained
from unfailing completion of the union. That system may be infinite,
and its normal forms need not be computable; the source uses their
existence for a mathematical completeness argument, not for the
decision algorithm. It also identifies layer-reduced forms, applying
collapse equations as far as possible, as an effective alternative.
The chapter expressly sketches this argument and refers to the
1996 primary for the full completion details, and to
Schmidt--Schau\ss/Kirchner--Ringeissen for layer reduction.
Those exact constructions are not identified with the tower below.
\cite[pp.~518--519]{BaaderSnyder2001}
\end{hbmethod}

\begin{hbdefinition}{def:p12-comb-tower}{An explicit equational free-product construction}
For the completeness proof, take a countably infinite generator set Z.
Write \(F_i(S)=T(\Sigma_i,S)/{=_{E_i}}\) for the pure free algebra on
a nonempty set S of named generators. Generators here are algebraic
atoms; they are not assumed to be nullary symbols in \(\Sigma_i\).
Define increasing sets of nodes of each colour i as follows:
\[
 S_i^0=Z,\qquad
 N_i^0=F_i(Z)\setminus Z,\qquad
 S_i^{k+1}=Z\cup\bigcup_{j\ne i}N_j^k,\qquad
 N_i^{k+1}=F_i(S_i^{k+1})\setminus S_i^{k+1}.
\]
In this display a free algebra's named generators are identified with
their generator classes. Nongenerator classes are tagged with their
colour, and older nodes are identified with their images under the
natural free-algebra inclusions. The next lemma proves that these
inclusions are injective and that old nongenerators remain nongenerators;
the displayed construction does not assume these facts.

Put \(S_i^\infty=Z\cup\bigcup_{j\ne i}\bigcup_kN_j^k\) and
\(D=Z\cup\bigcup_i\bigcup_kN_i^k\). There is a bijection
\(e_i:D\longrightarrow F_i(S_i^\infty)\): a node of colour i is
its pure class; any other node and any z in Z is a generator class.
Transport the \(\Sigma_i\)-operations along \(e_i\) to D.
Write \(\nu(t)\) for evaluation of a mixed term t in D under the
generator assignment. A generator has rank 0; a node first appearing
in \(N_i^k\) has rank \(k+1\). Each node admits a pure i-term
representative with foreign-colour atoms of strictly smaller rank.
The ranks are mathematical proof data, not an asserted effective
normalization algorithm.
\end{hbdefinition}

\begin{hblemma}{thm:p12-comb-free-product}{proved=annex:p12.comb.free-product}{The tower is the combined free algebra, including collapsing theories}
The tower construction is well defined. With its transported operations,
D is free on Z in the variety of E. In particular,
\(\nu(s)=\nu(t)\) iff \(s=_Et\), each reduct is the indicated pure
free algebra, and every element has a finite mixed-term representative.
Pure abstraction to distinct atoms reflects \(E_i\)-equality. These
statements remain valid for irregular and collapsing identities; no
minimum-support or syntactic-root preservation assumption is used.
\end{hblemma}

\begin{hbmethod}{def:p12-comb-witness}{Constructing the successful branch from an arbitrary solution}
Let \(\tau\) solve the purified problem; its finite image variables
are included in Z. Evaluate its images in the proved free algebra D.
Partition shared variables by equality of these values. Label a
representative by its value's colour; values in Z receive label 1.
Choose an order extending increasing rank, breaking ties arbitrarily.

For component i, each foreign shared value is a distinct generator
of \(F_i(S_i^\infty)\). Name it by its fixed representative constant.
Choose a pure representative of each own shared value at its first
tower level. For nonshared variables choose any pure representative
of \(e_i(\nu(x\tau))\). Replace every remaining atom in these finitely
many representatives by a fresh local parameter, consistently and
injectively, and retain the foreign shared constants. The resulting
substitution \(\theta_i\) is a local LCR solution: the next lemma proves
both the equations and every restriction. In particular, a written
root of an unnormalized target image is not used as a signature label.
\end{hbmethod}

\begin{hblemma}{thm:p12-comb-projection}{proved=annex:p12.comb.projection}{Actual normalized projection and the constant restrictions}
Every combined solution \(\tau\) yields the branch and local solutions
\(\theta_i\) just constructed. The abstraction solves
\(\Gamma_i\sigma_\Pi\) modulo \(E_i\), and no forbidden foreign shared
constant appears. Hence combined solvability implies local solvability
in at least one enumerated branch.
\end{hblemma}

\begin{hbtheorem}{thm:p12-comb-completeness}{proved=annex:p12.comb.completeness}{Complete local families give a complete combined family}
In each branch let \(C_i\) be complete for its component LCR problem,
and use disjoint fresh ranges as above. The union, over all branches,
of all ordered merges of tuples from \(C_1\times\cdots\times C_n\),
with identification and purification reconstruction, is complete modulo E
on the original problem variables. In particular, if all component LCR
problems have finite complete sets, every elementary E-problem has a
finite complete set. The proof supplies an actual global factor
substitution, not a claim that separate local factors are compatible.
\end{hbtheorem}

\begin{hbtheorem}{thm:p12-comb-decision}{proved=annex:p12.comb.decision}{Decision, finite generation, and the free-symbol extension}
If each component LCR decision procedure terminates and is correct,
the finite branch search decides elementary E-unifiability. If these
problems are in NP, so is elementary E-unifiability for this fixed
finite collection of theories. If component procedures terminate with
finite complete sets, finite combined-set generation terminates.
This implies at-most-finitary type, but the returned union need not
already be minimal. Under these terminating finite-set algorithms,
minimalization is also effective: to test whether one answer covers
another on the original scope, rename the first answer's range variables
fresh, freeze the second answer's range variables as distinct new
free constants, and decide the resulting finite image equations by
the same combination decision procedure. Quotient by mutual coverage
and retain only the most-general classes of the finite output.
This comparison is a new finite solvability call, not a requirement
to solve the word problem by computing the infinite tower.

For a finite input containing additional uninterpreted function symbols,
add the empty theory over precisely those symbols as one disjoint
component. Its LCR solver computes a syntactic MGU, freshens its range,
and checks forbidden-constant occurrences. A failed check excludes every
syntactic unifier; a successful MGU is a complete singleton LCR family.
Thus the same construction handles general E-unification under the
same component LCR hypotheses. This use of the syntactic interface
does not replace any equational component solver.
\cite[Proposition~6.1 and Theorems~6.2--6.3, pp.~516--517]{BaaderSnyder2001}
\end{hbtheorem}

\begin{hbtheorem}{thm:p12-comb-ac-lcr}{proved=annex:p12.comb.ac-lcr}{The positive pure-AC construction supplies the missing LCR interface}
For a pure unit-free AC component with designated free constants,
run the finite algorithm of Theorem~\ref{thm:p12-ac-coverage}, then
discard each answer containing a forbidden constant in an input
variable's image. This yields a finite complete AC-LCR family and
terminates. Thus the proved combination construction computes finite
complete sets for arbitrary finite nested terms over finitely many
pairwise distinct AC symbols and additional free symbols/constants.
No empty AC term or unit is introduced. This supplies a new explicit
algorithm for that general domain; it does not identify it with
Stickel's historical recursive reconciliation algorithm.
\end{hbtheorem}

\begin{hbexample}{def:p12-comb-cycle}{Identification is not enough: the source's cyclic failure}
Let f be associative with no unit and g be free unary. The source
problem is
\[
g(f(y,y))=^{?}g(x),\qquad g(x)=^{?}g(y),\qquad x=^{?}f(y,y).
\]
Purification gives the f-block \(z=^{?}f(y,y),\ x=^{?}f(y,y)\)
and the g-block \(g(z)=^{?}g(x),\ g(x)=^{?}g(y)\).
Label x,z with f and y with g. Distinct fixed x,z make the g-block fail.
Identifying z with x removes that failure, but the local substitutions
\(x\mapsto f(y,y)\) and \(y\mapsto x\) form a cycle. If \(x<y\),
the f-substitution violates LCR; if \(y<x\), the g-substitution does.
Neither ordering admits this tuple. Indeed g's injectivity forces
\(x=_Ey\), while the f-equation would require an associative word to
equal its own double, impossible for a nonempty word. This is a genuine
unsatisfiable case, not a reason to discard the ordering stage.
\cite[pp.~514--515]{BaaderSnyder2001}
\end{hbexample}

\begin{hbexample}{def:p12-comb-success}{A successful merge with a nested alien term}
Combine unit-free AC on f with free unary g and free constant a.
For \(f(x,g(a))=^{?}f(g(a),g(y))\), one purification has
\(u=^{?}g(a),\ v=^{?}g(y)\) in the free block and
\(f(x,u)=^{?}f(u,v)\) in the AC block. Take the singleton partition
of u,v, label both free, and choose \(u<v\). The free answer is
\(u\mapsto g(a),\ v\mapsto g(p),\ y\mapsto p\).
The AC answer, treating u,v as distinct constants, is \(x\mapsto v\).
The ordered merge is
\(u\mapsto g(a),\ v\mapsto g(p),\ x\mapsto g(p),\ y\mapsto p\).
Its reconstruction solves the original equation by AC. Cancellation
of the common g(a) bag atom shows that every original solution has
\(x=_{AC}g(y)\), so this displayed answer also covers every solution
of this particular problem. The preceding theorem, not this example,
establishes completeness for all finite nested AC inputs.
\end{hbexample}

\begin{hbexample}{def:p12-comb-collapse}{Why root labels and empty-support shortcuts are unsafe}
Take \(E_1=\{f(x)\approx x\}\) and free g in \(\Sigma_2\).
The written image \(f(g(z))\) has root colour 1 but denotes the same
combined element as \(g(z)\), of colour 2. Tower evaluation records that
collapse before assigning an owner. As an irregular example,
\(E_1=\{h(x)\approx h(y)\}\) has a nongenerator class represented
by h(z), independent of which generator z is chosen, even if there
are no nullary symbols at all. The tower keeps this class as a node;
it does not demand a ground representative or assume equality of
subalgebra intersections over an empty set. Both theories are
nontrivial and are admitted by the construction.
\end{hbexample}

\begin{hbexample}{def:p12-comb-ag-filter-failure}{The AC filter is not a generic equational LCR solver}
Use additive AG notation, with identity 0, inverse -, and a fixed free
constant c. The equation \(x+c=^{?}y\) has the ordinary AG-MGU
\(x\mapsto p,\ y\mapsto p+c\). With order \(y<c<x\), its written
y-image contains the forbidden constant c, whereas x may contain c.
The instance \(p\mapsto -c\) gives the permitted unifier
\(x\mapsto -c,\ y\mapsto0\). Thus discarding this ordinary AG-MGU loses a permitted
target. A correct AG-LCR solver must construct a different permitted
family; fixed-atom persistence cannot be assumed for signed counts.
Here 0 denotes the group identity e, not an extra signature constant.
\end{hbexample}

\begin{proofboundary}{def:p12-comb-boundaries}{What the combination theorem does and does not supply}
The full proofs here establish the chapter's elementary/general
disjoint-theory combination results through an explicit equational
free-product route. They do not identify this tower with the exact
1992/1996 ordered-completion construction or reproduce the full
quasi-free-structure algorithm of the 1998 primary. Predicate constraints,
disunification, shared interpreted signatures, and the source's
optimization/complexity examples require their own arguments.
Specifically, the source reports polynomial ACUI-LCR decision versus
NP-complete general ACUI decision to demonstrate an inherent cost of
combination. It reports a polynomial deterministic method for regular,
collapse-free theories with polynomial MGU computation, and practical
optimizations by Kepser--Richts. These hypotheses and source claims
are retained, but their complexity lower bounds and optimizing
algorithms are not proved by the finite-branch NP argument above.
\cite[pp.~516--517]{BaaderSnyder2001}

The AC-LCR bridge above uses the proved persistence of fixed atoms in
nonnegative AC multiplicities. It cannot silently be reused for AG:
signed occurrences can cancel after parameter instantiation. It supplies
neither the other monoidal source theorems, arbitrary-theory hereditary
transport, nor a historical manifestation comparison. Finite branch search
does not make a nonterminating component enumeration terminate, and
the mathematical tower is not an algorithm for the word problem in
arbitrary undecidable theories.
\end{proofboundary}


\begin{proofboundary}{def:p12-eu-method-boundaries}{Independent methods and remaining source-presentation duties}
The semiring construction uses the category over \(\operatorname{Sig}(E)\),
not an arbitrary free-function extension, which need not remain
semi-additive. Source \S6's combination method assumes nontrivial theories
over pairwise disjoint signatures and component solvers with linear
constant restrictions. Abstraction of alien subterms is followed by a
partition of shared variables, theory labels for representatives, and
a linear ordering excluding cyclic image dependencies. Separate
component solvability does not ensure a compatible combined substitution.
The actual purification, ordered merge and global factorization in
\cref{thm:p12-comb-purification,thm:p12-comb-soundness,
thm:p12-comb-completeness} now supply an independent full equational
free-product proof. With the proved component AC-LCR and syntactic
interfaces, \cref{thm:p12-comb-ac-lcr} supplies finite complete
families for arbitrary finite nested/multiple AC symbols and free
symbols; \cref{thm:p12-comb-decision} gives effective minimalization
under its terminating local-algorithm hypotheses.

The chapter's original normalized-root/strict-subterm argument remains
a source sketch, separately recorded in
\cref{def:p12-comb-source-witness}. Its exact 1992/1996 ordered-
completion and layer-reduction presentations are not identified with
the proved tower. Stickel's historical recursive-reconciliation
algorithm and the distinct 1981 journal presentation are not replaced
by this independent route. Signed/group LCR, other monoidal source
constructions, nondisjoint and predicate/disunification extensions,
and the source's optimizing/complexity arguments retain the precise
boundaries of \cref{def:p12-comb-boundaries}.
\cite[\S6, pp.~513--519]{BaaderSnyder2001}

Source \S4 also distinguishes its arbitrary-theory transformation
calculus, narrowing over ground-convergent presentations, and basic
narrowing with retained skeleton positions and an accumulator.
Enumeration completeness does not imply termination. The actual G
lift and directed B/N lift are proved in
\cref{thm:p12-eu-g-completeness,thm:p12-eu-bn-completeness}; neither
the syntactic MGU theorem nor the semiring proof would supply them
by itself. The arbitrary-theory hereditary transports remain distinct
from the supplied convergent-regime treatment in
\cref{sec:p12-eu-hereditary-source}. Source Corollary~4.14's
basic-to-ordinary passage has the exact same-family state/occurrence
transport of \cref{thm:p12-en-trace-transport}; it preserves a given
trace's matching-first scheduling and does not extend the convergence
domain. The weaker original-ground directed claim is refuted by
\cref{def:p12-eu-ground-domain-failure}. Initial open reduced witnesses
are supplied under the actual source ordering in
\cref{thm:p12-eb-initial-normalization}, and need no fresh constants.
The remaining raw inverse-basic requirement is global conversion,
protected-history and progress for E-equal different open normal
endpoints, recorded in \cref{def:p12-eb-progress-boundary}.
The left-linear regular countercase in
\cref{thm:p12-eb-progress-countercase} exhibits a successful
normal-witness inverse branch that increases the source M; it
does not refute raw inverse-basic completeness itself.

Terminal answer normalization in
\cref{thm:p12-en-terminal-normalization} preserves equational
coverage under its image-reduction hypotheses; it is not replacement
of literal active S-bindings or identification of primitive N with
a historical normalization-based algorithm. Finite matching blocks
do not imply search termination, as the infinite minimal family in
\cref{thm:p12-en-infinite-search} shows. The additional strategy
refinements and exact ordered closure retain their separately stated
dependencies.
\cite[\S4, pp.~488--503]{BaaderSnyder2001}
\end{proofboundary}
